% \iffalse meta-comment
% !TEX program  = XeLaTeX
%<*internal>
\iffalse
%</internal>
%<*readme>
xpinyin
=======

`xpinyin` is a LaTeX package written to simplify the input of Hanyu Pinyin.
The package provides macros to automatically add pinyin to Chinese characters.
It can only be used in conjunction with `xeCJK` or `CJK/CJKutf8` package.

Given that the implementation of pinyin package of `CJK` bundle is not very
well, it seems that xpinyin is a good replacement of it.

You can read the package manual (in Chinese) for more detailed explanations.

Contributing
------------

This package is a part of the [ctex-kit](https://github.com/CTeX-org/ctex-kit) project.

Issues and pull requests are welcome.

Copyright and Licence
---------------------

    Copyright (C) 2012-2022 by Qing Lee <sobenlee@gmail.com>
    ----------------------------------------------------------------------

    This work may be distributed and/or modified under the
    conditions of the LaTeX Project Public License, either
    version 1.3c of this license or (at your option) any later
    version. This version of this license is in
       http://www.latex-project.org/lppl/lppl-1-3c.txt
    and the latest version of this license is in
       http://www.latex-project.org/lppl.txt
    and version 1.3 or later is part of all distributions of
    LaTeX version 2005/12/01 or later.

    This work has the LPPL maintenance status "maintained".

    The Current Maintainer of this work is Qing Lee.

    This package consists of the file  xpinyin.dtx,
                 and the derived files xpinyin.sty,
                                       xpinyin.pdf,
                                       xpinyin.ins,
                                       xpinyin.lua,
                                       xpinyin.db,
                                       xpinyin-query.db,
                                       xpinyin-database.def,
                                       xpinyin-query.def and
                                       README.md (this file).

%</readme>
%<*internal>
\fi
\begingroup
  \def\temp{LaTeX2e}
\expandafter\endgroup\ifx\temp\fmtname\else
\csname fi\endcsname
%</internal>
%<*install>

%<install>\begingroup\expandafter\expandafter\expandafter\endgroup
%<install>\expandafter\ifx\csname directlua\endcsname\relax
%<install>  \begingroup
%<install>  \newlinechar=10 %
%<install>  \errmessage{^^J
%<install>    ********************************************^^J
%<install>    * LuaTeX is required to unpack this package.^^J
%<install>    ********************************************^^J
%<install>  }
%<install>  \endgroup
%<install>  \csname @@end\expandafter\endcsname\expandafter\end
%<install>\fi

\input ctxdocstrip %

\let\MetaPrefix\relax

\preamble

    Copyright (C) 2012-2022 by Qing Lee <sobenlee@gmail.com>
--------------------------------------------------------------------------

    This work may be distributed and/or modified under the
    conditions of the LaTeX Project Public License, either
    version 1.3c of this license or (at your option) any later
    version. This version of this license is in
       http://www.latex-project.org/lppl/lppl-1-3c.txt
    and the latest version of this license is in
       http://www.latex-project.org/lppl.txt
    and version 1.3 or later is part of all distributions of
    LaTeX version 2005/12/01 or later.

    This work has the LPPL maintenance status "maintained".

    The Current Maintainer of this work is Qing Lee.

--------------------------------------------------------------------------

\endpreamble

\postamble

    This package consists of the file  xpinyin.dtx,
                 and the derived files xpinyin.sty,
                                       xpinyin.pdf,
                                       xpinyin.ins,
                                       xpinyin.lua,
                                       xpinyin.db,
                                       xpinyin-query.db,
                                       xpinyin-database.def,
                                       xpinyin-query.def and
                                       README.md.
\endpostamble

\declarepostamble\emptypostamble
\endpostamble

\def\MetaPrefix{-- }

\generate
  {
    \usedir{source/latex/xpinyin}
    \usepostamble\emptypostamble
    \file{xpinyin.lua} {\from{\jobname.dtx}{lua}}
  }

\let\MetaPrefix\DoubleperCent

%<install>\directlua
%<install>  {
%<install>    if not kpse.find_file("xpinyin.db") then
%<install>      dofile(kpse.find_file("xpinyin.lua"))
%<install>    end
%<install>  }
%<install>\generate
%<install>  {
%<install>    \usedir{tex/latex/xpinyin}
%<install>    \usepostamble\emptypostamble
%<install>    \file{xpinyin-database.def}
%<install>      {
%<install>        \from{\jobname.dtx} {database}
%<install>        \from{xpinyin.db}   {}
%<install>      }
%<install>    \file{xpinyin-query.def}
%<install>      {
%<install>        \from{\jobname.dtx}      {querydb}
%<install>        \from{xpinyin-query.db}  {}
%<install>      }
%<install>  }

\generate
  {
%</install>
%<*internal>
    \usedir{source/latex/xpinyin}
    \file{xpinyin.ins} {\from{\jobname.dtx}{install}}
%</internal>
%<*install>
    \usedir{tex/latex/xpinyin}
    \file{xpinyin.sty} {\from{\jobname.dtx}{package}}
    \nopreamble\nopostamble
    \usedir{doc/latex/xpinyin}
    \file{README.md}   {\from{\jobname.dtx}{readme}}
  }

\endbatchfile
%</install>
%<*internal>
\fi
%</internal>
%<package>\NeedsTeXFormat{LaTeX2e}
%<package>\RequirePackage{expl3}
%<+package>\GetIdInfo$Id$
%<package>  {Automatically add pinyin to Chinese characters}
%<package>\ProvidesExplPackage{\ExplFileName}
%<package>  {\ExplFileDate}{3.1}{\ExplFileDescription}
%<database>\ProvidesFile{xpinyin-database.def}%
%<database>  [2022/07/14 v3.1 xpinyin database]
%<querydb>\ProvidesFile{xpinyin-query.def}%
%<querydb>  [2026/08/06 v3.1 xpinyin query database]
%<*driver>
\documentclass{ctxdoc}
%% 手册里「查询汉字的读音」一节的示例要用查询命令, 所以这里加 query 选项.
\usepackage[query]{xpinyin}
\newfontfamily\PinYinFont{TeX Gyre Adventor}
\xpinyinsetup{font=\PinYinFont,multiple=\color{red}}
%% 手册里多处要排带声调的 ü（ǖ ǘ ǚ ǜ, U+01D6/01D8/01DA/01DC）——它们正是
%% 「ü 的字母部分记作 v」这条约定要说明的对象。ctxdoc 的等宽字体 cmun 没有
%% 这四个字形, 缺字形时 XeTeX 静默丢掉分音符, 恰好把那几处想展示的区别抹掉
%% (实测手册里出现十余条 Missing character 警告)。
%%
%% 改用一个有这些字形的等宽字体。只在本手册生效, 不影响 ctxdoc 的其他使用者;
%% 手册里 xelatex 不支持 luaotfload 的 fallback 机制, 所以直接换字体而不是加后备。
\setmonofont{DejaVuSansMono}[
  Extension      = .ttf,
  UprightFont    = *,
  BoldFont       = *-Bold,
  ItalicFont     = *-Oblique,
  BoldItalicFont = *-BoldOblique,
  Scale          = MatchLowercase,
  HyphenChar     = None]
\begin{document}
  \DocInput{\jobname.dtx}
  \IndexLayout
  \PrintIndex
\end{document}
%</driver>
% \fi
%
% \GetFileId{xpinyin.sty}
%
% \title{\bfseries\pkg{xpinyin} 宏包}
% \author{李清\\ \path{sobenlee@gmail.com}}
% \date{\filedate\qquad\fileversion\thanks{\ctexkitrev{\ExplFileVersion}.}}
% \maketitle
%
% \begin{documentation}
%
% \section{简介}
%
% \pkg{xpinyin} 是一个 \LaTeX 宏包，提供了为汉字自动注音的功能。
%
% \section{基本用法}
%
% \pkg{xpinyin} 支持采用 \texttt{GBK} 和 \texttt{UTF-8} 编码的 \TeX 源文件，建议
% 总是使用 \texttt{UTF-8}。如果使用 \LaTeX 或 \pdfLaTeX\ 的编译方式，则根据编码的
% 情况，\pkg{xpinyin} 依赖 \package{CJK} 或者 \package[cjk]{CJKutf8} 宏包。
% 如果使用 \XeLaTeX，则依赖 \package{xeCJK} 宏包。如果它们没有在 \pkg{xpinyin}
% 之前被载入，\pkg{xpinyin} 将根据编译方式自动选择，\LaTeX 或 \pdfLaTeX\ 将使用
% \pkg{CJKutf8}。
%
% \pkg{xpinyin} 还依赖 \package{l3kernel} 和 \package{l3packages}，使用 (pdf)\LaTeX
% 下的 \texttt{GBK} 编码时，还将依赖 \package{xCJK2uni}。
%
% 需要注意的是，\pkg{xpinyin} 缺省将拼音的字体设置为与文档的主字体（\tn{normalfont}）相同，
% 所以为了保证声调字母的正确输出，应该选用合适的西文主字体。也可以通过将在下一节介绍的
% \meta{font} 选项来单独设置拼音的字体。
%
% \begin{minipage}[t]{\dimexpr(\linewidth-\parindent-2em)/2\relax}
% \XeLaTeX 下的简单示例：
% \begin{verbatim}[frame=single,gobble=3]
%   \documentclass{article}
%   \usepackage{xeCJK}
%   \usepackage{xpinyin}
%   \setmainfont{CMU Serif}
%   \setCJKmainfont{SimSun}
%
%   \begin{document}
%   \xpinyin*{汉语拼音示例}
%   \end{document}
% \end{verbatim}
% \end{minipage}\qquad
% \begin{minipage}[t]{\dimexpr(\linewidth-\parindent-2em)/2\relax}
% (pdf)\LaTeX 下的简单示例：
% \begin{verbatim}[frame=single,gobble=3]
%   \documentclass{article}
%   \usepackage{CJKutf8}
%   \usepackage{xpinyin}
%   \usepackage[T1]{fontenc}
%   \usepackage{lmodern}
%
%   \begin{document}
%   \begin{CJK}{UTF8}{gbsn}
%   \xpinyin*{汉语拼音示例}
%   \end{CJK}
%   \end{document}
% \end{verbatim}
% \hrule height \dp\strutbox width 0pt \relax
% \end{minipage}
%
% 运行上述示例要求系统安装了设置的字体，源文件用 |UTF-8| 编码保存，使用相应的编译方式。
%
% \pkg{xpinyin} 可以与 \package{ctex} 宏包或文档类共同使用，使用方式与上面类似。
%
% \section{用户手册}
%
% \begin{function}{pinyinscope}
%   \begin{syntax}
%     \tn{begin}\{pinyinscope\}\oarg{options}
%     .....
%     \tn{end}\{pinyinscope\}
%   \end{syntax}
%   为 \env{pinyinscope} 环境中的汉字自动注音。例如
%   \begin{Example}[frame=single,numbers=left,gobble=5]
%     \begin{pinyinscope}
%     列位看官：你道此书从何而来？说起根由，虽近荒唐，细按则深有趣味。
%     待在下将此来历注明，方使阅者\xpinyin{了}{liao3}然不惑。
%     \end{pinyinscope}
%   \end{Example}
% \end{function}
%
%   可选项 \meta{options} 用于局部设置拼音的格式，将在下面说明。
%
% \begin{function}{\xpinyin}
%   \begin{syntax}
%     \tn{xpinyin} \oarg{options} \Arg{单个汉字} \Arg{拼音}
%     \tn{xpinyin*} \oarg{options} \Arg{文字}
%   \end{syntax}
%   对于多音字，可以使用 \tn{xpinyin} 为其设置拼音；而 \tn{xpinyin*} 相当于
%   \env{pinyinscope} 环境的命令形式。\tn{xpinyin} 可以在 \env{pinyinscope} 环境和
%   \tn{xpinyin*} 中使用。例如，
%   \begin{SideBySideExample}
%     \xpinyin{长}{chang2}\\
%     \xpinyin*{甄士隐梦幻识通灵}\\
%     \xpinyin*{\xpinyin{重}{zhong4}要}
%   \end{SideBySideExample}
% \end{function}
%
% \begin{function}{\pinyin}
%   \begin{syntax}
%     \tn{pinyin} \oarg{options} \Arg{拼音}
%   \end{syntax}
%   用于输出拼音，为了输入的方便 \texttt{\"u} 可以用 |v| 代替。例如，
%   \begin{SideBySideExample}
%     \pinyin{lv2zi}\\
%     \pinyin{nv3hai2zi}
%   \end{SideBySideExample}
% \end{function}
%
% \begin{function}{\setpinyin}
%   \begin{syntax}
%     \tn{setpinyin} \Arg{汉字} \Arg{拼音}
%   \end{syntax}
%   \pkg{xpinyin} 宏包的拼音数据（\file{xpinyin-database.def}）来源于 \texttt{Unicode}
%   的 \texttt{Unihan} 数据库\footnotemark 中的 \file{Unihan_Readings.txt} 文件。对于多
%   音字，一般来说这个文件选用的是常用读音。可以使用 \tn{setpinyin} 来设置多音字的首选读音。
% \end{function}
%
% \footnotetext{\url{http://www.unicode.org/Public/UNIDATA/Unihan.zip}}
%
% \begin{function}{\xpinyinsetup}
%   \begin{syntax}
%     \tn{xpinyinsetup} \{ <key_1>=<val_1>, <key_2>=<val_2>, ... \}
%   \end{syntax}
%   用于在导言区或文档中，设置拼音的格式。目前可以设置的 \meta{key} 如下介绍。
% \end{function}
%
% \begin{function}{ratio}
%   \begin{syntax}
%     ratio = \Arg{number}
%   \end{syntax}
%   设置拼音字体大小与当前正文字体大小的比例，缺省值是 |0.4|。
% \end{function}
%
% \begin{function}{vsep}
%   \begin{syntax}
%     vsep = \Arg{dimen}
%   \end{syntax}
%   设置拼音的基线与汉字基线的间距，缺省值是 |1 em|。
% \end{function}
%
% \begin{function}{hsep}
%   \begin{syntax}
%     hsep = \Arg{skip}
%   \end{syntax}
%   设置注音汉字之间的间距，缺省值与 \tn{CJKglue} 的值相同。为了断行时行末的对齐，设置
%   的 \meta{skip} 最后有一定的弹性。例如
%   \begin{Example}[frame=single,numbers=left,gobble=5]
%     \xpinyin*[ratio={.7},hsep={.5em plus .1em},vsep={1.1em}]{贾雨村风尘怀闺秀}
%   \end{Example}
% \end{function}
%
% \begin{function}{pysep}
%   \begin{syntax}
%     pysep = \Arg{glue}
%   \end{syntax}
%   设置 \tn{pinyin} 输出的相邻两个汉语拼音的空白，缺省值是一个空格。
% \end{function}
%
% \begin{function}{font}
%   \begin{syntax}
%     font = \Arg{font}
%   \end{syntax}
%   设置拼音的字体，缺省值是 \tn{normalfont}，即以正文西文字体相同。为了保证拼音能正确
%   输出，最好选用收字量较大的西文字体。
% \end{function}
%
% \begin{function}{format}
%   \begin{syntax}
%     format = \Arg{format}
%   \end{syntax}
%   设置拼音的其他格式，例如颜色等，缺省值为空。
% \end{function}
%
% \begin{function}{multiple}
%   \begin{syntax}
%     multiple = \Arg{format}
%   \end{syntax}
%   设置多音字拼音的其他格式，缺省值为空。可以通过这个选项来提醒校正多音字的拼音。
%   注意这个选项会覆盖 \opt{format} 设置的格式。例如本文档设置多音字拼音的颜色是
%   红色（需要载入 \pkg{color} 宏包）：
%   \begin{frameverb}
%   \xpinyinsetup{multiple={\color{red}}}
%   \end{frameverb}
% \end{function}
%
% \begin{function}[added = 2014-01-12]
%   {footnote}
%   \begin{syntax}
%     footnote = <\TFF>
%   \end{syntax}
%   是否对拼音环境中的脚注（\tn{footnote}）汉字加上拼音。缺省值为 \texttt{false}。
%   更一般的情况，请使用 \tn{disablepinyin}。
% \end{function}
%
% \changes{v3.2}{2026/07/13}{为 \tn{disablepinyin} 增加 star 版本：
%   用于将其作用域覆盖到 \tn{xpinyin} 自身。（\#265）}
%
% \begin{function}[added = 2014-01-12, updated = 2026-07-13]
%   {\disablepinyin, \enablepinyin}
%   \tn{disablepinyin} 和 \tn{enablepinyin} 分别用于在拼音环境（\env{pinyinscope}）
%   中临时取消和恢复对汉字的注音。
%   \tn{disablepinyin} 还额外提供带星号的形式，可将禁用范围扩展到 \tn{xpinyin} 宏自身。
% \end{function}
%
% \changes{v3.2}{2026/08/06}{新增查询汉字读音的四个命令
%   \tn{xpinyinvalue}、\tn{xpinyininitial}、\tn{xpinyinshengmu} 与
%   \tn{xpinyinyunmu}，可用于中文索引的分组与排序；数据表由 \opt{query}
%   选项按需载入。（\#550）}
%
% \section{查询汉字的读音}
%
% 上面几节的命令都是把拼音排在汉字上方。如果只想知道某个汉字\emph{读什么}，
% 而不想排出注音，可以用本节的四个查询命令。它们返回一段文字，不排版任何别的东西，
% 因此可以写在其他命令的参数里。用于索引排序键时有一点要注意，见“给索引条目加拼音排序键”一节。
%
% 这组命令需要一份额外的数据表，加载宏包时用 \opt{query} 选项要求载入：
% \begin{quote}
%   \verb|\usepackage[query]{xpinyin}|
% \end{quote}
% 不写这个选项时，注音功能照常可用，但调用查询命令会报错。
% 之所以要分开，是因为这份表有一兆多字节，只用注音功能的文档不必为它付出载入时间。
%
% \pkg{xpinyin} 本身支持 \XeTeX{} 与 \pdfTeX{} 两条路线，而查询命令只在 \XeTeX{} 下提供。
% \pdfTeX{} 加 \pkg{CJKutf8} 的路线上，一个汉字由多个字节组成，
% 取不到本组命令所需的码位，此时使用 \opt{query} 选项会报错。
%
% （\LuaTeX{} 目前不在本宏包的支持范围内，加载时就会中止，与查询命令无关。）
%
% \begin{function}[added = 2026-08-06]{\xpinyinvalue}
%   \begin{syntax}
%     \tn{xpinyinvalue} \Arg{单个汉字}
%     \tn{xpinyinvalue*} \Arg{单个汉字}
%   \end{syntax}
%   返回这个汉字的拼音。带星号的形式返回全部读音，用逗号分隔。
%   \begin{SideBySideExample}
%     中 的读音是 \xpinyinvalue{中}。\\
%     行 有几个读音：\xpinyinvalue*{行}。
%   \end{SideBySideExample}
% \end{function}
%
% \begin{function}[added = 2026-08-06]{\xpinyininitial}
%   \begin{syntax}
%     \tn{xpinyininitial} \Arg{单个汉字}
%     \tn{xpinyininitial*} \Arg{单个汉字}
%   \end{syntax}
%   返回拼音的第一个字母。带星号的形式返回全部读音各自的首字母，
%   重复的只保留一个。
%
%   这个命令适合给索引条目自动加分组前缀：全部汉字的首字母都落在
%   \texttt{a} 到 \texttt{z} 之内，因此可以据此把中文条目分成若干组。
%   \begin{SideBySideExample}
%     汉 归入 \xpinyininitial{汉} 组，\\
%     安 归入 \xpinyininitial{安} 组。\\
%     行 跨越 \xpinyininitial*{行} 两组。
%   \end{SideBySideExample}
% \end{function}
%
% \begin{function}[added = 2026-08-06]{\xpinyinshengmu, \xpinyinyunmu}
%   \begin{syntax}
%     \tn{xpinyinshengmu} \Arg{单个汉字}
%     \tn{xpinyinyunmu} \Arg{单个汉字}
%   \end{syntax}
%   分别返回声母与韵母，两者都提供带星号的形式。
%   零声母的字（“安”“儿”一类）声母为空。
%   \begin{SideBySideExample}
%     中：声母 \xpinyinshengmu{中}，
%     韵母 \xpinyinyunmu{中}。\\
%     安：声母为空，
%     韵母 \xpinyinyunmu{安}。
%   \end{SideBySideExample}
% \end{function}
%
% \begin{function}[added = 2026-08-06]{tone}
%   \begin{syntax}
%     tone = mark|number|none|v
%   \end{syntax}
%   声调用什么形式表示，缺省值为 \opt{mark}。四种取值的区别如下：
%
%   \begin{center}
%   \begin{tabular}{llll}
%     \hline
%     取值 & 含义 & 中 & 女 \\
%     \hline
%     \opt{mark}   & 调号标在字母上 & \texttt{zhōng}  & \texttt{nǚ}  \\
%     \opt{number} & 调号写成数字   & \texttt{zhong1} & \texttt{nv3} \\
%     \opt{none}／\opt{v} & 不带声调  & \texttt{zhong}  & \texttt{nv}  \\
%     \hline
%   \end{tabular}
%   \end{center}
%
%   \opt{none} 与 \opt{v} 目前结果相同。保留两个名字是因为“不带声调”和
%   “把 \texttt{ü} 写成 \texttt{v}”本是两件事，只是本宏包内部一律用 \texttt{v}
%   记录 \texttt{ü}，于是两者一致。
%
%   \opt{number} 形式里 \texttt{ü} 写作 \texttt{v}，与 \tn{pinyin} 的输入约定一致，
%   所以查询结果可以直接交回 \tn{pinyin}：
%   \begin{SideBySideExample}
%     \xpinyinsetup{tone=number}
%     \xpinyinvalue{女} 交回 \tn{pinyin}
%     得到 \pinyin{nv3}。
%   \end{SideBySideExample}
% \end{function}
%
% \begin{function}[added = 2026-08-06]{scheme}
%   \begin{syntax}
%     scheme = literal|official
%   \end{syntax}
%   声母与韵母按什么规则切分，缺省值为 \opt{literal}。这个选项只影响
%   \tn{xpinyinshengmu} 与 \tn{xpinyinyunmu}，不影响另外两个命令。
%
%   两种规则的差别只在 \texttt{y} 和 \texttt{w}。\opt{literal} 把它们当作声母，
%   切分结果拼回去就是原来的拼音。\opt{official} 依据《汉语拼音方案》，
%   把 \texttt{y} 和 \texttt{w} 视为介音 \texttt{i}、\texttt{u}、\texttt{ü}
%   的书写形式，声母为空，并把韵母还原成音位形式。
%
%   \begin{center}
%   \begin{tabular}{lllll}
%     \hline
%     & & \multicolumn{2}{c}{\opt{literal}} & \opt{official} \\
%     汉字 & 拼音 & 声母 & 韵母 & 韵母 \\
%     \hline
%     王 & \texttt{wáng} & \texttt{w} & \texttt{áng} & \texttt{uang} \\
%     月 & \texttt{yuè}  & \texttt{y} & \texttt{uè}  & \texttt{ve}   \\
%     鱼 & \texttt{yú}   & \texttt{y} & \texttt{ú}   & \texttt{v}    \\
%     安 & \texttt{ān}   &            & \texttt{ān}  & \texttt{an}   \\
%     \hline
%   \end{tabular}
%   \end{center}
%
%   \opt{official} 模式下的韵母\emph{不带声调}。还原出来的形式可能以
%   \texttt{v} 开头（\texttt{yue} 还原成 \texttt{ve}），而调号只能标在元音上，
%   \tn{pinyin} 也只在 \texttt{l} 与 \texttt{n} 之后才把 \texttt{v} 当作
%   \texttt{ü}，所以给这种形式标调会排出错的字形。同样的原因，
%   \opt{literal} 模式下以 \texttt{v} 开头的韵母（“女”的韵母）也不带声调。
%
%   \texttt{yo} 一类音节按“\texttt{y} 后接其他元音则 \texttt{y} 变 \texttt{i}”
%   这条通则处理成 \texttt{io}。这个音节的实际读法在语言学上有不同意见，
%   本宏包不对它单独取值。
% \end{function}
%
% \begin{function}[added = 2026-08-06]{query}
%   \begin{syntax}
%     query
%   \end{syntax}
%   宏包选项，要求载入查询表。只能在 \tn{usepackage} 时给出，
%   不能中途用 \tn{xpinyinsetup} 打开：数据表在加载宏包时读入一次。
% \end{function}
%
% \subsection{关于可展开性}
%
% 查询命令可以放在需要参数展开成字面文字的位置，例如 \pkg{bib2gls} 的字段。
% 除了下面两种情况，四个命令都是完全可展开的：
%
% \begin{itemize}
%   \item \opt{tone} 取 \opt{mark} 时。带调号的输出要经过 \tn{pinyin} 的标调过程，
%     它把内容排版出来而不是返回文字。
%   \item 带星号的形式。它需要去掉重复项，用到内部的列表变量。
% \end{itemize}
%
% 也就是说，需要字面文字时请选 \opt{number}、\opt{none} 或 \opt{v}，
% 并使用不带星号的形式。
%
% \subsection{给索引条目加拼音排序键}
%
% \tn{index} \emph{不会}展开参数。它先用 \tn{@sanitize} 把参数变成字符，
% 再原样写进 \file{.idx}，所以直接写
% \begin{quote}
%   \verb|\index{\xpinyinvalue{汉}@汉字}|
% \end{quote}
% 得到的是 \verb|\indexentry{\xpinyinvalue{汉}@汉字}{1}|——排序键里是这串命令本身，
% 而不是 \texttt{han}。这与查询命令可不可展开无关，是 \tn{index} 的既有行为。
%
% 要用拼音作排序键，得自己先展开一次再交给 \tn{index}：
% \begin{quote}
% \begin{verbatim}
% \newcommand\pyindex[2]{%
%   \begingroup
%     \xpinyinsetup{tone=none}%
%     \edef\pyindexkey{#1}%
%     \expandafter\index\expandafter{\pyindexkey @#2}%
%   \endgroup}
% ...
% 汉字\pyindex{\xpinyinvalue{汉}}{汉字}
% 安全\pyindex{\xpinyinvalue{安}}{安全}
% \end{verbatim}
% \end{quote}
% 这样写进 \file{.idx} 的是 \verb|\indexentry{han@汉字}{1}|，
% 经 \pkg{makeindex} 排序后「安全」「汉字」「中国」按 \texttt{an}、\texttt{han}、
% \texttt{zhong} 的顺序排列。
%
% 中间变量请用一个普通的名字。若写成 \tn{@tempa} 这类含 |@| 的名字，
% 正文里 |@| 不是字母，\verb|\edef\@tempa{...}| 定义到的并不是
% \tn{@tempa}，排序键会多出一个空格（\verb|han @汉字|）；
% \pkg{makeindex} 按字面比较，这会让它与手写的 \verb|han@汉字| 分成两个条目。
% 要用这类名字就得把整段放在 \tn{makeatletter} 与 \tn{makeatother} 之间。
%
% \subsection{多音字}
%
% 不带星号的形式取首选读音，与 \tn{xpinyin} 的取值一致；哪个读音算首选，
% 由 \file{Unihan} 数据库决定，也可以用 \tn{setpinyin} 改写。
%
% 需要注意的是，多音字的不同读音常常落在不同的首字母上。
% 数据库里的五千多个多音字中，约有三分之二属于这种情况。
% 因此按首字母给索引分组时，要留意条目可能应当出现在两处。
% 带星号的形式给出全部读音，便于自行判断：
% \begin{SideBySideExample}
%   行 的全部读音：\xpinyinvalue*{行}\\
%   涉及的首字母：\xpinyininitial*{行}
% \end{SideBySideExample}
%
% \end{documentation}
%
% \StopEventually{}
%
% \begin{implementation}
%
% \section{代码实现}
%
% \changes{v3.2}{2026/08/04}{补齐独立回归测试，并接入按 tag 构建发布包所需的
%   版本一致性校验。宏包代码本身没有变化。}
%
%    \begin{macrocode}
%<*package>
%<@@=xpinyin>
%    \end{macrocode}
%
%    \begin{macrocode}
\msg_new:nnn { xpinyin } { l3-too-old }
  {
    Support~package~'expl3'~too~old. \\\\
    Please~update~an~up~to~date~version~of~the~bundles\\\\
    'l3kernel'~and~'l3packages'\\\\
    using~your~TeX~package~manager~or~from~CTAN.
  }
\@ifpackagelater { expl3 } { 2019/03/05 } { }
  { \msg_error:nn { xpinyin } { l3-too-old } }
%    \end{macrocode}
%
%    \begin{macrocode}
\msg_new:nnn { xpinyin } { engine-not-supported }
  { Engine~`\c_sys_engine_str'~is~not~yet~supported,~xpinyin~will~abort! }
\bool_lazy_or:nnF
  { \sys_if_engine_xetex_p: }
  { \sys_if_engine_pdftex_p: }
  { \msg_critical:nn { xpinyin } { engine-not-supported } }
%    \end{macrocode}
%
% 查询命令需要额外的数据表，只有加载 \opt{query} 选项时才会读入。
% 未加载时直接报错，而不是让命令展开成空——后者会让排序键静默变成空串，
% 出错的地方离原因很远。
%    \begin{macrocode}
\msg_new:nnn { xpinyin } { query-not-loaded }
  {
    Pinyin~query~table~is~not~loaded. \\\\
    Load~the~package~with~the~`query'~option:\\\\
    \iow_char:N \\usepackage[query]{xpinyin}
  }
\msg_new:nnn { xpinyin } { query-no-reading }
  { No~pinyin~reading~for~`#1'~in~the~database. }
%    \end{macrocode}
% 可展开的报错在终端上只显示到行宽为止（实测 \verb|take exactly...| 之后被截断），
% 所以关键信息要放在第一行开头、且尽量短。完整说明留给第二行，
% 它在 \file{.log} 里仍是完整的。
%    \begin{macrocode}
\msg_new:nnn { xpinyin } { query-bad-argument }
  {
    Query~one~character~at~a~time~(got~`#1'). \\\\
    Pinyin~query~commands~take~exactly~one~character.
  }
%    \end{macrocode}
%
% 查询命令只在 \XeTeX{} 下提供。取汉字码位要用 |`#1|，
% 而 \pdfTeX{} 加 \pkg{CJKutf8} 的路线上一个汉字是多个字节，
% 这个写法拿不到码位；包内的 \cs{@@_to_unicode:n} 依赖 \cs{CJK@plane}，
% 需要字符级钩子提供的上下文，也不能在这里直接调用。
%    \begin{macrocode}
\msg_new:nnn { xpinyin } { query-needs-unicode }
  {
    Pinyin~query~commands~are~only~available~with~XeTeX. \\\\
    They~are~not~available~with~pdfTeX.
  }
%    \end{macrocode}
%
%    \begin{macrocode}
\cs_if_exist:NF \NewDocumentCommand
  { \RequirePackage { xparse } }
%    \end{macrocode}
%
% \begin{variable}{\c_@@_tone_prop}
% 重音标记 |\`|、|\'| 和 |\=| 在 \env{tabbing} 环境中被移作他用，为避免错误，
% 我们使用内部命令 \tn{@tabacckludge} 或 \tn{a} 来定义。
%    \begin{macrocode}
\prop_const_from_keyval:Nn \c_@@_tone_prop
  {
    { ā } = { \@tabacckludge= a } ,
    { á } = { \@tabacckludge' a } ,
    { ǎ } = { \v a } ,
    { à } = { \@tabacckludge` a } ,
    { ō } = { \@tabacckludge= o } ,
    { ó } = { \@tabacckludge' o } ,
    { ǒ } = { \v o } ,
    { ò } = { \@tabacckludge` o } ,
    { ē } = { \@tabacckludge= e } ,
    { é } = { \@tabacckludge' e } ,
    { ě } = { \v e } ,
    { è } = { \@tabacckludge` e } ,
    { ū } = { \@tabacckludge= u } ,
    { ú } = { \@tabacckludge' u } ,
    { ǔ } = { \v u } ,
    { ù } = { \@tabacckludge` u } ,
    { ḿ } = { \@tabacckludge' m } ,
    { ń } = { \@tabacckludge' n } ,
    { ň } = { \v n } ,
    { ǹ } = { \@tabacckludge` n } ,
    { ī } = { \@tabacckludge= { \i } } ,
    { í } = { \@tabacckludge' { \i } } ,
    { ǐ } = { \v { \i } } ,
    { ì } = { \@tabacckludge` { \i } } ,
    { ü } = { \" u } ,
    { ǖ } = { \@tabacckludge= { \" u } } ,
    { ǘ } = { \@tabacckludge' { \" u } } ,
    { ǚ } = { \v { \" u } } ,
    { ǜ } = { \@tabacckludge` { \" u } }
  }
%    \end{macrocode}
% \end{variable}
%
%
% \begin{macro}{\@@_UTF_char:nn}
%    \begin{macrocode}
\cs_new_protected:Npn \@@_UTF_char:nn #1#2
  {
    \cs_if_exist:cF { u8:#1 }
      { \tl_const:cn { u8:#1 } {#2} }
  }
%    \end{macrocode}
% \end{macro}
%
%
% \begin{macro}{\@@_GBK_char:nn}
%    \begin{macrocode}
\cs_new_protected:Npn \@@_GBK_char:nn #1#2
  {
    \@@_UTF_char:nn {#1} {#2}
    \exp_args:Nx \@@_GBK_char_aux:nn { \tl_head:n {#1} } {#1}
  }
\cs_new_protected:Npn \@@_GBK_char_aux:nn #1#2
  { \exp_args:Nf \@@_GBK_char_aux:nnn { \int_eval:n { `#1 } } {#1} {#2} }
\cs_new_protected:Npn \@@_GBK_char_aux:nnn #1#2#3
  {
    \cs_if_exist:cF { @@_UTF_ #1 :w }
      {
        \exp_args:Nf \@@_GBK_char_def:nnn
          {
            \int_case:nn { \tl_count:n {#3} }
              {
                { 2 } { ##1 }
                { 3 } { ##1##2 }
                { 4 } { ##1##2##3 }
              }
          }
          {#1} {#2}
        \exp_args:Nc \@@_save_UTF_cs:Nn { @@_UTF_ #1 :w } {#1}
        \tl_gput_right:Nx \c_@@_reset_UTF_catcode_tl
          { \char_set_catcode:nn {#1} { \char_value_catcode:n {#1} } }
        \char_set_catcode_active:n {#1}
      }
  }
\cs_new_protected:Npn \@@_GBK_char_def:nnn #1#2#3
  {
    \cs_new_protected:cpn { @@_UTF_ #2 :w } #1
      { \use:c { u8: \tl_to_str:n { #3#1 } } }
  }
\tl_new:N \c_@@_reset_UTF_catcode_tl
%    \end{macrocode}
% \end{macro}
%
% \begin{macro}{\@@_save_UTF_cs:Nn}
%    \begin{macrocode}
\group_begin:
\char_set_catcode_active:n { 126 }
\cs_new_protected:Npn \@@_save_UTF_cs:Nn #1#2
  {
    \group_begin:
    \char_set_lccode:nn { 126 } {#2}
    \tex_lowercase:D
      {
        \group_end:
        \tl_gput_right:Nn \c_@@_reset_UTF_cs_tl { \cs_set_eq:NN ~ #1 }
      }
  }
\group_end:
\tl_new:N \c_@@_reset_UTF_cs_tl
%    \end{macrocode}
% \end{macro}
%
%    \begin{macrocode}
\bool_new:N \g_@@_GBK_bool
\@ifpackageloaded { xeCJK }
  { \AtEndOfPackage { \@@_adjust_xeCJK_hook: } }
  {
    \@ifpackageloaded { CJKutf8 }
      {
        \prop_map_function:NN \c_@@_tone_prop \@@_UTF_char:nn
        \AtEndOfPackage { \@@_adjust_CJK_hook: }
      }
      {
        \@ifpackageloaded { CJK }
          {
            \RequirePackage { xCJK2uni }
            \prop_map_function:NN \c_@@_tone_prop \@@_GBK_char:nn
            \AtEndOfPackage
              {
                \tl_put_right:Nn \l_@@_pinyin_box_hook_tl
                  { \c_@@_reset_UTF_cs_tl }
                \@@_adjust_CJK_hook:
                \tl_use:N \c_@@_reset_UTF_catcode_tl
              }
            \bool_gset_true:N \g_@@_GBK_bool
          }
          {
            \sys_if_engine_xetex:TF
              {
                \RequirePackage { xeCJK }
                \AtEndOfPackage { \@@_adjust_xeCJK_hook: }
              }
              {
                \RequirePackage { CJKutf8 }
                \prop_map_function:NN \c_@@_tone_prop \@@_UTF_char:nn
                \AtEndOfPackage { \@@_adjust_CJK_hook: }
              }
          }
      }
  }
%    \end{macrocode}
%
% \begin{variable}{\l_@@_tmpa_box,\l_@@_tmpb_box}
%    \begin{macrocode}
\box_new:N \l_@@_tmpa_box
\box_new:N \l_@@_tmpb_box
%    \end{macrocode}
% \end{variable}
%
% \begin{macro}{\@@_width:Nn}
%    \begin{macrocode}
\cs_new_protected:Npn \@@_width:Nn #1#2
  {
    \hbox_set:Nn \l_@@_tmpa_box {#2}
    #1 = \box_wd:N \l_@@_tmpa_box
  }
%    \end{macrocode}
% \end{macro}
%
% \begin{macro}{\@@_make_pinyin_box:nnn}
%    \begin{macrocode}
\cs_new_protected:Npn \@@_make_pinyin_box:nnn #1#2#3
  {
    \@@_leavevmode:
    \hbox_overlap_right:n
      {
        \hbox_set:Nn \l_@@_tmpa_box
          { \@@_CJKsymbol_hook: \@@_save_CJKsymbol:n {#2} }
        \hbox_set:Nn \l_@@_tmpb_box
          {
            \l_@@_pinyin_box_hook_tl
            \@@_select_font:
            \l_@@_format_tl
            \clist_if_exist:cT { c_@@_multiple_ #1 _clist }
              { \l_@@_multiple_tl }
            {#3}
          }
        \dim_compare:nNnT
          { \box_wd:N \l_@@_tmpb_box } >
          { \box_wd:N \l_@@_tmpa_box + \l_@@_CJKglue_dim }
          {
            \box_resize_to_wd_and_ht:Nnn \l_@@_tmpb_box
              { \box_wd:N \l_@@_tmpa_box + \l_@@_CJKglue_dim }
              { \box_ht:N \l_@@_tmpb_box }
          }
        \box_move_up:nn { \l_@@_vsep_tl }
          {
            \hbox_to_wd:nn { \box_wd:N \l_@@_tmpa_box }
              { \tex_hss:D \box_use_drop:N \l_@@_tmpb_box \tex_hss:D }
          }
      }
  }
\tl_new:N \l_@@_pinyin_box_hook_tl
\sys_if_engine_pdftex:T
  {
    \tl_put_right:Nn \l_@@_pinyin_box_hook_tl
      { \cs_set_eq:NN \CJK@plane \tex_undefined:D }
  }
\cs_generate_variant:Nn \@@_make_pinyin_box:nnn { x }
%    \end{macrocode}
% \end{macro}
%
% \begin{macro}{\@@_CJKsymbol:n}
%    \begin{macrocode}
\cs_new_protected:Npn \@@_CJKsymbol:n #1
  { \@@_CJKsymbol:xn { \@@_to_unicode:n {#1} } {#1} }
\cs_new_protected:Npn \@@_CJKsymbol:nn #1#2
  {
    \@@_make_pinyin_box:nnn {#1} {#2} { \use:c { c_@@_ #1 _tl } }
    \@@_save_CJKsymbol:n {#2}
  }
\cs_generate_variant:Nn \@@_CJKsymbol:nn { x }
%    \end{macrocode}
% \end{macro}
%
% \begin{macro}{pinyinscope}
%    \begin{macrocode}
\NewDocumentEnvironment { pinyinscope } { O { } }
  {
    \keys_set:nn { xpinyin } {#1}
    \enablepinyin
  }
  { }
%    \end{macrocode}
% \end{macro}
%
% \begin{macro}{\xpinyin}
%    \begin{macrocode}
\NewDocumentCommand \xpinyin { s O { } m }
  {
    \mode_leave_vertical:
    \bool_if:NTF \l_@@_enable_all_bool
      {
        \IfBooleanTF {#1}
          {
            \group_begin:
              \keys_set:nn { xpinyin } {#2}
              \enablepinyin
              #3
            \group_end:
          }
          {
            \group_begin:
              \keys_set:nn { xpinyin } {#2}
              \bool_if:NF \l_@@_enable_bool
                { \@@_width:Nn \l_@@_CJKglue_dim { \CJKglue } }
              \@@_single_aux:nn {#3}
          }
      }
      { \IfBooleanTF {#1} {#3} { \use_i:nn {#3} } }
  }
%    \end{macrocode}
% \end{macro}
%
% \begin{variable}{\l_@@_enable_bool, \l_@@_enable_all_bool}
%    \begin{macrocode}
\bool_new:N \l_@@_enable_bool
\bool_new:N \l_@@_enable_all_bool
\bool_set_true:N \l_@@_enable_all_bool
%    \end{macrocode}
% \end{variable}
%
% \begin{macro}{\@@_CJKglue:}
%    \begin{macrocode}
\cs_new_protected:Npn \@@_CJKglue:
  { \skip_horizontal:n { \l_@@_hsep_tl } }
%    \end{macrocode}
% \end{macro}
%
% \begin{macro}{\enablepinyin}
%    \begin{macrocode}
\NewDocumentCommand \enablepinyin { }
  {
    \bool_if:NF \l_@@_enable_all_bool
      { \bool_set_true:N \l_@@_enable_all_bool }
    \bool_if:NF \l_@@_enable_bool
      {
        \tl_if_empty:NF \l_@@_hsep_tl
          {
            \cs_set_eq:NN \@@_save_CJKglue: \CJKglue
            \cs_set_eq:NN \CJKglue \@@_CJKglue:
          }
        \@@_width:Nn \l_@@_CJKglue_dim { \CJKglue }
        \@@_replace_CJKsymbol:
        \@@_restore_footnote:
        \bool_set_true:N \l_@@_enable_bool
      }
  }
%    \end{macrocode}
% \end{macro}
%
% \begin{macro}{\disablepinyin}
%    \begin{macrocode}
\NewDocumentCommand \disablepinyin { s }
  {
    \bool_if:NT #1
      { \bool_set_false:N \l_@@_enable_all_bool }
    \bool_if:NT \l_@@_enable_bool
      {
        \cs_if_eq:NNT \CJKglue \@@_CJKglue:
          { \cs_set_eq:NN \CJKglue \@@_save_CJKglue: }
        \@@_restore_CJKsymbol:
        \bool_set_false:N \l_@@_enable_bool
      }
  }
%    \end{macrocode}
% \end{macro}
%
% \begin{macro}{\@@_restore_footnote:}
%    \begin{macrocode}
\cs_new_protected:Npn \@@_restore_footnote:
  {
    \bool_if:NF \l_@@_footnote_bool
      { \tl_put_left:Nn \@parboxrestore { \l_@@_restore_footnote_tl } }
  }
%    \end{macrocode}
% \end{macro}
%
% \begin{variable}{\l_@@_restore_footnote_tl}
%    \begin{macrocode}
\tl_new:N \l_@@_restore_footnote_tl
\tl_set:Nn \l_@@_restore_footnote_tl
  {
    \int_compare:nNnT \tex_currentgrouptype:D = { 11 }
      { \disablepinyin }
  }
%    \end{macrocode}
% \end{variable}
%
% \begin{variable}{\l_@@_CJKglue_dim}
%    \begin{macrocode}
\dim_new:N \l_@@_CJKglue_dim
%    \end{macrocode}
% \end{variable}
%
% \begin{macro}{\@@_single_aux:nn}
%    \begin{macrocode}
\cs_new_protected:Npn \@@_single_aux:nn #1#2
  {
      \@@_replace_CJKsymbol_single:n {#2}
      #1
    \group_end:
  }
\cs_new_protected:Npn \@@_replace_CJKsymbol_single_aux:n #1
  {
    \bool_if:NF \l_@@_enable_bool { \@@_replace_CJKsymbol: }
    \cs_set_protected:Npn \CJKsymbol ##1
      { \@@_single_CJKsymbol:nn {##1} {#1} }
  }
\cs_new_protected:Npn \@@_single_CJKsymbol:nn #1#2
  {
    \@@_make_pinyin_box:xnn
      { \@@_to_unicode:n {#1} } {#1} { \@@_pinyin:n {#2} }
    \@@_save_CJKsymbol:n {#1}
  }
%    \end{macrocode}
% \end{macro}
%
% \begin{macro}{\@@_replace_CJKsymbol_aux:}
%    \begin{macrocode}
\cs_new_protected:Npn \@@_replace_CJKsymbol_aux:
  {
    \cs_set_eq:NN \@@_save_CJKsymbol:n \CJKsymbol
    \cs_set_eq:NN \CJKsymbol \@@_CJKsymbol:n
  }
%    \end{macrocode}
% \end{macro}
%
% \begin{macro}{\@@_restore_CJKsymbol_aux:}
%    \begin{macrocode}
\cs_new_protected:Npn \@@_restore_CJKsymbol_aux:
  { \cs_set_eq:NN \CJKsymbol \@@_save_CJKsymbol:n }
%    \end{macrocode}
% \end{macro}
%
% \begin{macro}{\@@_select_font_xetex:}
%    \begin{macrocode}
\cs_new_protected:Npn \@@_select_font_xetex:
  {
    \cs_if_exist_use:cF { \l_@@_coor_tl }
      {
        \tl_set:Nx \l_@@_current_coor_tl { \l_@@_coor_tl }
        \@@_select_font_aux:
        \int_compare:nNnF { \tex_XeTeXfonttype:D \tex_font:D } = \c_zero_int
          {
            \exp_last_unbraced:NNV
            \cs_gset_eq:cN \l_@@_current_coor_tl \tex_font:D
          }
      }
  }
%    \end{macrocode}
% \end{macro}
%
% \begin{macro}{\@@_select_font_aux:}
%    \begin{macrocode}
\cs_new_protected:Npn \@@_select_font_aux:
  {
    \fontsize
      { \l_@@_ratio_tl \tex_dimexpr:D \f@size pt \scan_stop: }
      { \f@baselineskip }
    \normalfont
    \l_@@_font_tl
    \selectfont
  }
%    \end{macrocode}
% \end{macro}
%
% \begin{macro}{\@@_to_unicode_xetex:n}
%    \begin{macrocode}
\cs_new:Npn \@@_to_unicode_xetex:n #1
  { \int_to_arabic:n { `#1 } }
%    \end{macrocode}
% \end{macro}
%
% \begin{macro}{\@@_UTF_to_unicode:n,\@@_UTFchar_to_unicode:n}
%    \begin{macrocode}
\cs_new:Npn \@@_UTF_to_unicode:n #1
  {
    \int_to_arabic:n
      { \exp_args:No \int_from_hex:n { \CJK@plane } * "100 + #1 }
  }
\cs_new:Npn \@@_UTFchar_to_unicode:n #1
  { \int_to_arabic:n { \@@_UTF_viii_to_unicode:NNNw #1 \q_stop } }
\cs_new:Npn \@@_UTF_viii_to_unicode:NNNw #1#2#3#4 \q_stop
  {
    \tl_if_empty:nTF {#4}
      { ( `#1 - "E0 ) * "1000 + ( `#2 - "80 ) * "40 + ( `#3 - "80 ) }
      { ( `#1 - "F0 ) * "4000 + ( `#2 - "80 ) * "1000 + ( `#3 - "80 ) * "40 + ( `#4 - "80 ) }
  }
%    \end{macrocode}
% \end{macro}
%
% \begin{macro}{\@@_GBK_to_unicode:n,\@@_GBKchar_to_unicode:n}
%    \begin{macrocode}
\cs_new:Npn \@@_GBK_to_unicode:n
  { \CJKtu_sfd_map:nn { \CJK@plane } }
\cs_new:Npn \@@_GBKchar_to_unicode:n
  { \CJKtu_char_to_unicode:n }
%    \end{macrocode}
% \end{macro}
%
% \begin{macro}{\@@_adjust_xeCJK_hook:}
%    \begin{macrocode}
\cs_new_protected:Npn \@@_adjust_xeCJK_hook:
  {
    \cs_new_eq:NN \@@_select_font:       \@@_select_font_xetex:
    \cs_new_eq:NN \@@_to_unicode:n       \@@_to_unicode_xetex:n
    \cs_new_eq:NN \@@_char_to_unicode:n  \@@_to_unicode:n
    \cs_new_eq:NN \@@_restore_CJKsymbol: \@@_restore_CJKsymbol_aux:
    \cs_new_eq:NN \@@_replace_CJKsymbol: \@@_replace_CJKsymbol_aux:
    \cs_new_eq:NN \@@_replace_CJKsymbol_single:n
                  \@@_replace_CJKsymbol_single_aux:n
    \tl_if_exist:NTF \l_xeCJK_current_font_tl
      {
        \tl_set:Nn \l_@@_coor_tl
          {
            ( \tl_to_str:N \l_@@_font_tl ) /
            \l_xeCJK_current_font_tl/\l_@@_ratio_tl
          }
      }
      {
        \tl_set:Nn \l_@@_coor_tl
          {
            ( \tl_to_str:N \l_@@_font_tl ) /
            \xeCJK@family/\f@series/\f@shape/\f@size/\l_@@_ratio_tl
          }
      }
    \cs_new_eq:NN \@@_leavevmode: \prg_do_nothing:
    \cs_new_protected:Npx \@@_CJKsymbol_hook:
      {
        \exp_not:N \makexeCJKinactive
        \exp_not:N \@@_reselect_CJK_font:
      }
  }
%    \end{macrocode}
% \end{macro}
%
% \changes{v3.2}{2026/08/05}{在 \pkg{xeCJK} 的 \texttt{AutoFallBack} 已切换到
%   后备字体时，量宽盒子不再重选主字体，修复注音被压缩到错误宽度的问题（\#997）。}
%
% \begin{macro}{\@@_reselect_CJK_font:}
% \cs{@@_CJKsymbol_hook:} 原本无条件调用 \tn{xeCJK@setfont}（或 \pkg{xeCJK} 的
% \texttt{\textbackslash xeCJK\_select\_font:}）重选 CJK 字体，为量宽盒子里的汉字
% 准备字体。问题在于后者的内部实现第一件事就是清除后备字体状态，会把
% \texttt{AutoFallBack} 刚切好的后备字体丢掉。于是启用 \texttt{AutoFallBack} 且当前
% 字符靠后备字体才有字形时，量宽盒子退回主字体、字形缺失，量出的是该字体的
% \emph{notdef} 宽度（实测 2.8pt，正确值 10.0pt；这个数值与标点无关——同一字体下
% 任何缺字形的字符都得 2.8pt，而真正的句点是 2.78pt），拼音随后被 \cs{box_resize_to_wd_and_ht:Nnn} 压缩到
% 这个错误宽度。注意可见的汉字本身不受影响：它由盒子外的
% \cs{@@_save_CJKsymbol:n} 输出，此时后备字体仍然有效。
%
% \texttt{\textbackslash xeCJK\_reset\_fallback\_font:} 是 \pkg{xeCJK} 用来标记
% 「当前处于后备字体状态」的开关：未启用后备字体时它等于 \cs{prg_do_nothing:}，
% \pkg{xeCJK} 切换到后备字体后它被重定义为「恢复该字体并清除标记」。据此判断：
% 已经在后备字体里就不重选，当前字体正是应该用来量宽的那一个；否则保持原有行为。
%
% 为什么保留「否则重选」这一支，而不是把重选整个删掉：实测在本 hook 运行的时点，
% 当前字体（\tn{fontname}\tn{font}）已经是 CJK 字体，因为 \pkg{xeCJK} 的
% interchar 机制进入 CJK 类时已经切好字体，本 hook 运行在那之后。已核对
% 十余种上下文（紧跟西文单词、标点、\tn{emph}、数学、\tn{textsf}、字号变化、
% \tn{mbox}／\tn{sbox}／\tn{hbox}／\tn{vbox}、\env{tabular}、\env{minipage}、
% \tn{section}、\env{pinyinscope}、嵌套注音、脚注、切换 CJK 字体族），
% \tn{fontname}\tn{font} 无一例外都已是 CJK 字体；把重选删掉时全部回归测试仍然通过。
% 也就是说这一支在当前的 \pkg{xeCJK} 行为下并无可观察效果，保留它是出于对
% \pkg{xeCJK} 内部行为变化的保守（\texttt{NFSS} 参数此时确实还是西文族，一旦
% \pkg{xeCJK} 改变切换时点，重选就重新变得必要），而不是因为已经观察到它必需。
% 相应地，回归测试无法覆盖「这一支被跳过」这种形态——记为已接受的覆盖缺口，
% 见 \file{llmdoc/reference/build-and-test.md} 的 \#997 一节。
%
% 这个开关在 \pkg{xeCJK} 里没有独立的文档条目，属于内部量，比
% \texttt{\textbackslash xeCJK\_select\_font:} 更容易在上游重构中改名。本包依赖的
% \pkg{xeCJK} 内部接口清单见 \file{MAINTAINING.md}；\pkg{xeCJK} 侧对「重选字体会
% 清除后备字体状态」这条机制边界的说明，见仓库内
% \file{llmdoc/architecture/xecjk-architecture.md} 的「后备字体 (Fallback)」一节。
%    \begin{macrocode}
\cs_new_protected:Npn \@@_reselect_CJK_font:
  {
    \cs_if_exist:NTF \xeCJK_reset_fallback_font:
      {
        \cs_if_eq:NNTF \xeCJK_reset_fallback_font: \prg_do_nothing:
          { \@@_select_CJK_font: }
          { }
      }
      { \@@_select_CJK_font: }
  }
\cs_new_protected:Npn \@@_select_CJK_font:
  {
    \cs_if_exist_use:NF \xeCJK_select_font:
      { \xeCJK@setfont }
  }
%    \end{macrocode}
% \end{macro}
%
% \begin{macro}{\@@_adjust_CJK_hook:}
%    \begin{macrocode}
\cs_new_protected:Npn \@@_adjust_CJK_hook:
  {
    \bool_if:NTF \g_@@_GBK_bool
      {
        \cs_new_eq:NN \@@_to_unicode:n      \@@_GBK_to_unicode:n
        \cs_new_eq:NN \@@_char_to_unicode:n \@@_GBKchar_to_unicode:n
      }
      {
        \cs_new_eq:NN \@@_to_unicode:n      \@@_UTF_to_unicode:n
        \cs_new_eq:NN \@@_char_to_unicode:n \@@_UTFchar_to_unicode:n
      }
    \cs_new_eq:NN \@@_select_font:    \@@_select_font_aux:
    \cs_new_eq:NN \@@_leavevmode:     \mode_leave_vertical:
    \cs_new_eq:NN \@@_CJKsymbol_hook: \prg_do_nothing:
    \@ifpackageloaded { CJKpunct }
      { \@@_adjust_CJKpunct_hook: }
      {
        \cs_new_eq:NN \@@_restore_CJKsymbol: \@@_restore_CJKsymbol_aux:
        \cs_new_eq:NN \@@_replace_CJKsymbol: \@@_replace_CJKsymbol_aux:
        \cs_new_eq:NN \@@_replace_CJKsymbol_single:n
                      \@@_replace_CJKsymbol_single_aux:n
        \AtBeginDocument
          {
            \@ifpackageloaded { CJKpunct }
              {
                \cs_undefine:N \@@_restore_CJKsymbol:
                \cs_undefine:N \@@_replace_CJKsymbol:
                \cs_undefine:N \@@_replace_CJKsymbol_single:n
                \@@_adjust_CJKpunct_hook:
              } { }
          }
      }
  }
%    \end{macrocode}
% \end{macro}
%
% \begin{macro}{\@@_adjust_CJKpunct_hook:}
%    \begin{macrocode}
\cs_new_protected:Npn \@@_adjust_CJKpunct_hook:
  {
    \cs_new_protected:Npn \@@_restore_CJKsymbol:
      {
        \int_compare:nNnTF { \CJKpunct@punctstyle } = { \CJKpunct@ps@plain }
          { \@@_restore_CJKsymbol_aux: }
          { \cs_set_eq:NN \CJKosymbol \@@_save_CJKsymbol:n }
      }
    \cs_new_protected:Npn \@@_replace_CJKsymbol:
      {
        \int_compare:nNnTF { \CJKpunct@punctstyle } = { \CJKpunct@ps@plain }
          { \@@_replace_CJKsymbol_aux: }
          {
            \cs_set_eq:NN \@@_save_CJKsymbol:n \CJKosymbol
            \cs_set_eq:NN \CJKosymbol \@@_CJKsymbol:n
          }
      }
    \cs_new_protected:Npn \@@_replace_CJKsymbol_single:n ##1
      {
        \int_compare:nNnTF { \CJKpunct@punctstyle } = { \CJKpunct@ps@plain }
          { \@@_replace_CJKsymbol_single_aux:n { ##1 } }
          {
            \bool_if:NF \l_@@_enable_bool
              { \cs_set_eq:NN \@@_save_CJKsymbol:n \CJKosymbol }
            \cs_set_protected:Npn \CJKosymbol ####1
              { \@@_single_CJKsymbol:nn { ####1 } { ##1 } }
          }
      }
  }
%    \end{macrocode}
% \end{macro}
%
% \begin{macro}{\pinyin}
%    \begin{macrocode}
\NewDocumentCommand \pinyin { O { } m }
  {
    \group_begin:
    \keys_set:nn { xpinyin } {#1}
    \l_@@_font_tl
    \l_@@_format_tl { }
    \selectfont
    \c_@@_reset_UTF_cs_tl
    \@@_pinyin:n {#2}
    \group_end:
  }
%    \end{macrocode}
% \end{macro}
%
% \begin{macro}{\@@_pinyin:n}
%    \begin{macrocode}
\cs_new_protected:Npn \@@_pinyin:n #1
  {
    \@@_pinyin_init:
    \bool_set_true:N \l_@@_first_bool
    \tl_set:Nn \l_@@_save_tl {#1}
    \@@_pinyin_aux:n #1 \q_recursion_tail \q_recursion_stop
  }
%    \end{macrocode}
% \end{macro}
%
% \begin{macro}{\@@_pinyin_aux:n}
%    \begin{macrocode}
\cs_new_protected:Npn \@@_pinyin_aux:n #1
  {
    \quark_if_recursion_tail_stop_do:nn {#1}
      {
        \bool_if:NTF \l_@@_first_bool { \l_@@_save_tl }
          { \tl_if_empty:NF \l_@@_item_tl { \l_@@_pysep_tl \l_@@_item_tl } }
      }
    \@@_if_number:nTF {#1}
      {
        \bool_if:NTF \l_@@_first_bool
          { \bool_set_false:N \l_@@_first_bool }
          { \l_@@_pysep_tl }
        \l_@@_pre_tl
        \@@_tone:Vn \l_@@_tone_tl {#1}
        \l_@@_post_tl
        \@@_pinyin_init:
      }
      {
        \int_compare:nNnTF
          { 0 \cs_if_exist_use:c { c_@@_ \tl_to_str:N \l_@@_tone_tl _tl } } >
          { 0 \cs_if_exist_use:c { c_@@_ \tl_to_str:n {#1} _tl } }
          { \tl_put_right:Nn \l_@@_post_tl {#1} }
          {
            \tl_set:Nn \l_@@_tone_tl {#1}
            \tl_set_eq:NN \l_@@_pre_tl \l_@@_item_tl
            \tl_clear:N \l_@@_post_tl
          }
        \tl_put_right:Nx \l_@@_item_tl { \@@_replace_v:n {#1} }
      }
    \@@_pinyin_aux:n
  }
%    \end{macrocode}
% \end{macro}
%
% \begin{macro}{\@@_tone:nn}
%    \begin{macrocode}
\cs_new_protected:Npn \@@_tone:nn #1#2
  { \use:c { @@_num_to_tone_ #1 :Nn } {#1} {#2} }
\cs_generate_variant:Nn \@@_tone:nn { V }
%    \end{macrocode}
% \end{macro}
%
% \begin{macro}{\@@_replace_v:n}
%    \begin{macrocode}
\cs_new:Npn \@@_replace_v:n #1
  {
    \str_if_eq:nnTF {#1} { v }
      {
        \str_case:onTF { \l_@@_item_tl }
          { { l } { } { n } { } { L } { } { N } { } }
          { \exp_not:n { ü } } { u }
      }
      { \exp_not:n {#1} }
  }
%    \end{macrocode}
% \end{macro}
%
% \begin{macro}{\@@_pinyin_init:}
%    \begin{macrocode}
\cs_new:Npn \@@_pinyin_init:
  {
    \tl_clear:N \l_@@_pre_tl   \tl_clear:N \l_@@_post_tl
    \tl_clear:N \l_@@_item_tl  \tl_clear:N \l_@@_tone_tl
  }
%    \end{macrocode}
% \end{macro}
%
% \begin{macro}{\@@_if_number:nTF}
%    \begin{macrocode}
\prg_new_conditional:Npnn \@@_if_number:n #1 { TF }
  {
    \if_int_compare:w \c_one_int < 1 \tl_to_str:n {#1} \exp_stop_f:
      \prg_return_true: \else: \prg_return_false: \fi:
  }
%    \end{macrocode}
% \end{macro}
%
% \begin{variable}{\l_@@_first_bool}
%    \begin{macrocode}
\bool_new:N \l_@@_first_bool
%    \end{macrocode}
% \end{variable}
%
% \begin{variable}
% {\c_@@_a_tl,\c_@@_o_tl,\c_@@_e_tl,\c_@@_i_tl,\c_@@_u_tl,\c_@@_v_tl}
%    \begin{macrocode}
\tl_const:Nn \c_@@_a_tl { 3 }
\tl_const:Nn \c_@@_o_tl { 2 }
\tl_const:Nn \c_@@_e_tl { 2 }
\tl_const:Nn \c_@@_i_tl { 1 }
\tl_const:Nn \c_@@_u_tl { 1 }
\tl_const:Nn \c_@@_v_tl { 1 }
%    \end{macrocode}
% \end{variable}
%
% \begin{macro}{\@@_num_to_tone:Nn}
%    \begin{macrocode}
\cs_new_protected:Npn \@@_num_to_tone:Nn #1#2
  {
    \if_case:w \int_eval:n { #2 - 1 } \exp_stop_f:
      \= {#1}  \or: \'{#1}  \or: \v {#1}  \or: \` {#1}  \else: #1 \fi:
  }
\tl_map_inline:nn { a o e u }
  { \cs_new_eq:cN { @@_num_to_tone_ #1 :Nn } \@@_num_to_tone:Nn }
\cs_new:Npn \@@_num_to_tone_i:Nn #1#2
  {
    \if_case:w \int_eval:n { #2 - 1 } \exp_stop_f:
      ī  \or: í  \or: ǐ  \or: ì \else: i \fi:
  }
\cs_new_protected:Npn \@@_num_to_tone_v:Nn #1#2
  {
    \str_case:onTF { \l_@@_pre_tl }
      { { l } { }  { n } { }  { L } { } { N } { } }
      {
        \if_case:w \int_eval:n { #2 - 1 } \exp_stop_f:
          ǖ  \or: ǘ  \or: ǚ  \or: ǜ \else: ü \fi:
      }
      { \@@_num_to_tone:Nn u {#2} }
  }
%    \end{macrocode}
% \end{macro}
%
% \begin{macro}{\xpinyinsetup}
%    \begin{macrocode}
\NewDocumentCommand \xpinyinsetup { m } { \keys_set:nn { xpinyin } {#1} }
%    \end{macrocode}
% \end{macro}
%
% \begin{macro}{ratio,vsep,hsep,pysep,font,format,multiple,footnote}
%    \begin{macrocode}
\clist_map_inline:nn
  { ratio , vsep , hsep , pysep , font , format , multiple }
  { \keys_define:nn { xpinyin } { #1 .tl_set:c = { l_@@_ #1 _tl } } }
\keys_define:nn { xpinyin }
  { footnote .bool_set:N = \l_@@_footnote_bool }
\keys_set:nn { xpinyin }
  {
    ratio   = .4 ,
    vsep    = 1 em ,
    pysep   = \c_space_tl ,
    font    = \normalfont ,
  }
%    \end{macrocode}
% \end{macro}
%
% \begin{macro}{query}
% 是否载入查询表。这是包选项，不能在文档中途改变：数据表只在加载宏包时读一次。
%    \begin{macrocode}
\keys_define:nn { xpinyin }
  { query .bool_gset:N = \g_@@_query_bool }
%    \end{macrocode}
% \end{macro}
%
% \begin{macro}[int]{\xpinyin_customary:nnn,\xpinyin_multiple:nnn}
%    \begin{macrocode}
\cs_new_protected:Npn \xpinyin_customary:nnn #1#2
  { \cs_gset_nopar:cpn { c_@@_ #2 _tl } }
\cs_new_protected:Npn \xpinyin_multiple:nnn #1#2
  { \cs_gset_nopar:cpn { c_@@_multiple_ #2 _clist } }
%    \end{macrocode}
% \end{macro}
%
%    \begin{macrocode}
\group_begin:
  \cs_set_eq:NN \XPYU  \xpinyin_customary:nnn
  \cs_set_eq:NN \XPYUM \xpinyin_multiple:nnn
  \file_input:n { xpinyin-database.def }
\group_end:
%    \end{macrocode}
%
% \begin{macro}[int]{\xpinyin_query:nn,\xpinyin_query_multiple:nn}
% 载入查询表时用的两个入口，与 \cs{xpinyin_customary:nnn} 的作用相同，
% 只是存的是“字母加声调数字”的形式，且不再需要汉字本身作参数。
%    \begin{macrocode}
\cs_new_protected:Npn \xpinyin_query:nn #1
  { \cs_gset_nopar:cpn { c_@@_query_ #1 _tl } }
\cs_new_protected:Npn \xpinyin_query_multiple:nn #1
  { \cs_gset_nopar:cpn { c_@@_query_multiple_ #1 _clist } }
%    \end{macrocode}
% \end{macro}
%
% \begin{macro}{\setpinyin}
%    \begin{macrocode}
\NewDocumentCommand \setpinyin { m m }
  {
    \tl_set:cn
      { c_@@_ \@@_char_to_unicode:n {#1} _tl }
      { \@@_pinyin:n {#2} }
  }
%    \end{macrocode}
% \end{macro}
%
% \subsection{查询命令}
%
% 下面四个命令回答“这个汉字读什么”，不排版任何内容。它们\emph{完全可展开}，
% 因此可以写进 \tn{csname} 的构造式、\tn{edef} 或 \pkg{bib2gls} 的字段——
% issue 550 提出这组命令，正是为了给中文索引算排序键。
%
% 注意 \tn{index} 本身\emph{不}展开参数（它先 \tn{@sanitize} 再原样写进
% \file{.idx}），所以要用拼音作排序键必须自己先 \tn{edef} 一次；
% 手册的“给索引条目加拼音排序键”一节给出了可用的写法。
%
% 可展开的代价是所有加工都必须在生成数据库时完成：去掉声调要把预组合字符
% 拆开，而 \pkg{expl3} 里做这件事的 \cs{char_to_nfd:N} 无法在可展开的上下文里
% 使用（实测嵌套在 \tn{edef} 里会直接出错）。因此查询表存的是
% “字母加声调数字”的形式（|zhong1|、|nv3|），四种输出都由它派生，
% 运行时只做取值与字符串处理。
%
% \begin{variable}{\l_@@_query_scheme_tl,\l_@@_query_tone_tl}
% 两个选项的当前值。\opt{scheme} 决定声母与韵母怎么切分，
% \opt{tone} 决定声调以什么形式出现。
%    \begin{macrocode}
\tl_new:N \l_@@_query_scheme_tl
\tl_new:N \l_@@_query_tone_tl
\keys_define:nn { xpinyin }
  {
%    \end{macrocode}
% 分支代码里用 |#1| 而不是 \cs{l_keys_choice_tl}：前者在选项解析时就已替换成
% 字面分支名，不依赖运行时的局部变量（见 \pkg{ctex} 的 \opt{space} 选项踩过的坑）。
%    \begin{macrocode}
    scheme .choices:nn =
      { literal , official }
      { \tl_set:Nn \l_@@_query_scheme_tl {#1} } ,
    tone   .choices:nn =
      { mark , number , none , v }
      { \tl_set:Nn \l_@@_query_tone_tl {#1} } ,
  }
\keys_set:nn { xpinyin } { scheme = literal , tone = mark }
%    \end{macrocode}
% \end{variable}
%
% \begin{macro}[int]{\@@_query_get:n,\@@_query_get_multiple:n}
% 取出某个汉字的“字母加声调数字”形式。参数是汉字本身。
% 查询表未载入时给出提示，而不是静默展开成空。
%    \begin{macrocode}
\cs_new:Npn \@@_query_get:n #1
  { \use:c { c_@@_query_ \int_eval:n { `#1 } _tl } }
\cs_new:Npn \@@_query_get_multiple:n #1
  { \use:c { c_@@_query_multiple_ \int_eval:n { `#1 } _clist } }
%    \end{macrocode}
% \end{macro}
%
% \begin{macro}[int]{\@@_query_last:n,\@@_query_base:n,\@@_query_base:nn,
%   \@@_query_tone:n,\@@_query_tone:nn}
% 把“字母加声调数字”拆成字母部分与声调数字。判断末字符是不是数字时比较码位，
% 而不是用 \cs{str_if_in:nnTF}——它确实不可展开，会让整条链失去可展开性。
% 数字 |1| 到 |4| 的码位都小于字母 |a| 的码位，所以一次比较就够。
% 轻声没有调号，这时字母部分就是全串，声调部分为空。
%    \begin{macrocode}
\cs_new:Npn \@@_query_last:n #1 { \str_range:nnn {#1} {-1} {-1} }
\cs_new:Npn \@@_query_base:n #1
  { \exp_args:Nf \@@_query_base:nn { \@@_query_last:n {#1} } {#1} }
\cs_new:Npn \@@_query_base:nn #1#2
  {
    \int_compare:nNnTF { `#1 } < { `5 }
      { \str_range:nnn {#2} {1} {-2} } {#2}
  }
\cs_new:Npn \@@_query_tone:n #1
  { \exp_args:Nf \@@_query_tone:nn { \@@_query_last:n {#1} } {#1} }
\cs_new:Npn \@@_query_tone:nn #1#2
  { \int_compare:nNnTF { `#1 } < { `5 } {#1} { } }
%    \end{macrocode}
% \end{macro}
%
% \begin{macro}[int]{\@@_query_mark:n}
% 把“字母加声调数字”还原成带调号的形式。声调标在哪个字母上有固定规则：
% |a| 和 |e| 优先，其次是 |ou| 里的 |o|，否则标在最后一个元音上。
% 这里复用包内已有的 \cs{@@_pinyin:n}，它正是为 \tn{pinyin} 处理同样的输入而写的，
% 于是标调规则只有一份实现。注意这一支不可展开。
%
% \cs{@@_pinyin:n} 期望的是完整音节。\opt{official} 模式切出来的韵母是音位形式，
% 可能以 |v| 开头（|yue| 还原成 |ve|），而 \tn{pinyin} 把开头的 |v| 当成
% 声母位置上的字母，排出来是 |uè| 而不是 |üè|。这种组合本身不是合法音节，
% 所以带调号的输出对它没有意义：\opt{official} 模式下的韵母改用无调形式，
% 并在手册里说明这一限制。
%    \begin{macrocode}
\cs_new_protected:Npn \@@_query_mark:n #1 { \@@_pinyin:n {#1} }
%    \end{macrocode}
% \end{macro}
%
% \begin{macro}[int]{\@@_query_shengmu:nn,\@@_query_yunmu:nn}
% 声母与韵母的切分。两个模式的区别只在 |y| 和 |w| 算不算声母，
% 所以用两张表：\texttt{sm1l} 收 |y|、|w|，\texttt{sm1o} 不收。
% |zh|、|ch|、|sh| 是两个字母的声母，必须先于单字母匹配。
%
% 用“控制序列存不存在”来查表，是因为 \cs{cs_if_exist_use:cF} 可展开，
% \cs{str_case:nnF} 也可展开，这里用控制序列查表是因为两字母声母与单字母声母
% 要分两级匹配，用“存不存在”表达更直接。
%    \begin{macrocode}
\clist_map_inline:nn { zh , ch , sh }
  { \cs_new:cpn { @@_query_sm2_ #1 : } { #1 } }
\clist_map_inline:nn { b,p,m,f,d,t,n,l,g,k,h,j,q,x,r,z,c,s }
  {
    \cs_new:cpn { @@_query_sm1literal_ #1 : }  { #1 }
    \cs_new:cpn { @@_query_sm1official_ #1 : } { #1 }
  }
\clist_map_inline:nn { y , w }
  { \cs_new:cpn { @@_query_sm1literal_ #1 : } { #1 } }
\cs_new:Npn \@@_query_shengmu:nn #1#2
  { \exp_args:Nnf \@@_query_shengmu_aux:nn {#1} { \@@_query_base:n {#2} } }
%    \end{macrocode}
% 有几个音节整个就是一个鼻音：|m|、|n|、|ng| 以及 |hm|、|hng|
% （“嗯”“哼”一类，数据库里共 11 个字）。它们没有声母与韵母之分，
% 整体作韵母处理，声母为空——否则按字母切分会把“嗯”的 |ng| 切成
% |n| 加 |g|，而 |g| 不是韵母（实测早先输出 \verb|sm=[n] ym=[g4]|）。
%
% 这一步同时消掉了另一个缺陷：这些音节从前被切掉声母后韵母为空串，
% 与声调数字拼接后只剩一个裸数字（U+5463 的韵母得到 |2|，
% \opt{tone} 为 \opt{mark} 时页面上排出孤立的“2”）。
% 整体作韵母之后不再出现空韵母，因此韵母那一侧不必再加判据
% （核对过全部 57493 个读音，按现在的切分规则无一为空）。
%    \begin{macrocode}
\cs_new:Npn \@@_query_shengmu_aux:nn #1#2
  {
    \cs_if_exist:cTF { @@_query_syllabic_ #2 : }
      { }
      {
        \cs_if_exist_use:cF { @@_query_sm2_ \str_range:nnn {#2} {1} {2} : }
          {
            \cs_if_exist_use:cF
              { @@_query_sm1 #1 _ \str_range:nnn {#2} {1} {1} : }
              { }
          }
      }
  }
\clist_map_inline:nn { m , n , ng , hm , hng }
  { \cs_new:cpn { @@_query_syllabic_ #1 : } { #1 } }
%    \end{macrocode}
% 韵母是去掉声母之后剩下的部分。\opt{official} 模式下若声母为空，
% 还要把 |y| 或 |w| 的写法还原成介音。
%    \begin{macrocode}
\cs_new:Npn \@@_query_yunmu:nn #1#2
  {
    \exp_args:Nnff \@@_query_yunmu_aux:nnn {#1}
      { \@@_query_shengmu:nn {#1} {#2} } { \@@_query_base:n {#2} }
  }
\cs_new:Npn \@@_query_yunmu_aux:nnn #1#2#3
  {
    \str_if_eq:eeTF {#1} { official }
      {
        \tl_if_empty:nTF {#2}
          { \@@_query_glide:n {#3} }
          { \str_range:nnn {#3} { \tl_count:n {#2} + 1 } {-1} }
      }
      { \str_range:nnn {#3} { \tl_count:n {#2} + 1 } {-1} }
  }
%    \end{macrocode}
% \end{macro}
%
% \begin{macro}[int]{\@@_query_glide:n}
% \opt{official} 模式下把 |y| 与 |w| 的写法还原成介音，依据《汉语拼音方案》：
% |y| 后接 |i| 时去掉 |y|；|y| 后接 |u| 时该 |u| 实为 |ü|，按包内约定写作 |v|；
% |y| 后接其他元音时 |y| 变 |i|；|w| 后接 |u| 时去掉 |w|，接其他元音时 |w| 变 |u|。
% 不以 |y| 或 |w| 开头的（|an|、|er| 一类）原样返回。
%
% |yo| 一类按“|y| 后接其他元音”这一条处理成 |io|。这个音节的实际读法在
% 语言学上有不同意见，本包按上面的通则统一处理，不对它单独取值。
%    \begin{macrocode}
\cs_new:cpn { @@_query_glide_y_i: } { i }
\cs_new:cpn { @@_query_glide_y_u: } { v }
\cs_new:cpn { @@_query_glide_w_u: } { u }
\cs_new:cpn { @@_query_medial_y: }  { i }
\cs_new:cpn { @@_query_medial_w: }  { u }
\cs_new:Npn \@@_query_glide:n #1
  {
    \exp_args:Nff \@@_query_glide_aux:nn
      { \str_range:nnn {#1} {1} {1} } { \str_range:nnn {#1} {2} {-1} }
  }
\cs_new:Npn \@@_query_glide_aux:nn #1#2
  {
    \cs_if_exist:cTF { @@_query_medial_ #1 : }
      {
        \cs_if_exist:cTF
          { @@_query_glide_ #1 _ \str_range:nnn {#2} {1} {1} : }
          {
            \use:c { @@_query_glide_ #1 _ \str_range:nnn {#2} {1} {1} : }
            \str_range:nnn {#2} {2} {-1}
          }
          { \use:c { @@_query_medial_ #1 : } #2 }
      }
      { #1 #2 }
  }
%    \end{macrocode}
% \end{macro}
%
% \begin{macro}[int]{\@@_query_tone_apply:n}
% 按 \opt{tone} 的取值决定输出形式。\opt{number} 直接返回查询表里的原值；
% \opt{none} 去掉声调数字并保留 |v|；\opt{v} 与 \opt{none} 相同，
% 保留两个名字是因为“无声调”和“把 \texttt{ü} 写成 \texttt{v}”是两件事，
% 只是在本包的存储形式下结果一致；\opt{mark} 交给 \cs{@@_query_mark:n} 标调。
%    \begin{macrocode}
\cs_new:Npn \@@_query_tone_apply:n #1
  {
    \str_if_eq:eeTF { \l_@@_query_tone_tl } { number }
      { #1 }
      {
        \str_if_eq:eeTF { \l_@@_query_tone_tl } { mark }
          { \@@_query_mark_if_possible:n {#1} }
          { \@@_query_base:n {#1} }
      }
  }
%    \end{macrocode}
% \end{macro}
%
% \begin{macro}[int]{\@@_query_mark_if_possible:n,\@@_query_markable:n,
%   \@@_query_if_markable:nTF}
% 标调号要有可标的元音字母。整个音节就是一个鼻音时（|m|、|n|、|ng| 一类，
% 见前面的 \cs{@@_query_syllabic_} 表）没有元音，\cs{@@_pinyin:n} 找不到
% 落点，会把声调数字原样排出——\verb|\xpinyinvalue{嗯}| 得到 |ng4|，
% 页面上那个 |4| 就是个裸数字。
%
% 这与声母“永远不带声调”是同一件事的另一侧：判据都是“有没有可标调的元音”。
% 因此这里先检查，没有元音就退回无调形式（|ng| 而不是 |ng4|）。
%    \begin{macrocode}
\cs_new:Npn \@@_query_mark_if_possible:n #1
  {
    \@@_query_if_markable:nTF {#1}
      { \@@_query_mark:n {#1} }
      { \@@_query_base:n {#1} }
  }
%    \end{macrocode}
% 逐字符扫一遍看有没有元音字母。这里不能用 \cs{tl_if_in:nnTF} 或
% \cs{str_remove_all:nn}——两者都不可展开（实测前者在 \tn{edef} 里原样漏出
% 整条命令，后者把参数连花括号一起返回，长度反而变长），而查询命令必须保持
% 可展开。可展开的做法是用 \cs{str_range:nnn} 取单个字符，交给
% \cs{str_case:enF} 比对。
%    \begin{macrocode}
\cs_new:Npn \@@_query_markable:n #1
  { \@@_query_markable_aux:nn {#1} { 1 } }
\cs_new:Npn \@@_query_markable_aux:nn #1#2
  {
    \int_compare:nNnTF {#2} > { \str_count:n {#1} }
      { 0 }
      {
        \str_case:enF { \str_range:nnn {#1} {#2} {#2} }
          { {a}{1} {o}{1} {e}{1} {i}{1} {u}{1} {v}{1} }
          { \exp_args:Nnf \@@_query_markable_aux:nn {#1}
              { \int_eval:n { #2 + 1 } } }
      }
  }
\prg_new_conditional:Npnn \@@_query_if_markable:n #1 { TF }
  {
    \int_compare:nNnTF { \@@_query_markable:n {#1} } = { 1 }
      { \prg_return_true: } { \prg_return_false: }
  }
%    \end{macrocode}
% \end{macro}
%
% \begin{macro}[int]{\@@_query_part:nn}
% 按要取的部分（整串、首字母、声母、韵母）加工一个读音。
% |#1| 是部分的名字，|#2| 是查询表里的“字母加声调数字”。
%
% 首字母取的是字母部分的第一个字符，因此不受 \opt{tone} 影响；
% 其余三个都要经过 \cs{@@_query_tone_apply:n}。
%    \begin{macrocode}
\cs_new:Npn \@@_query_part:nn #1#2
  {
    \str_if_eq:nnTF {#1} { initial }
      { \exp_args:Nf \@@_query_head:n { \@@_query_base:n {#2} } }
      {
%    \end{macrocode}
% 声母只由辅音组成，标不了调号，所以\emph{永远}不带声调，四种 \opt{tone} 取值下
% 都一样，这一支不经过 \cs{@@_query_tone_apply:n}。
%
% 这一点是实测发现的：早先把声调一律接在切出来的片段后面，结果“中”的声母排成
% |zh1|——\cs{@@_pinyin:n} 找不到可以标调的元音，就把数字原样输出了。
%    \begin{macrocode}
        \str_if_eq:nnTF {#1} { shengmu }
          { \exp_args:NV \@@_query_shengmu:nn \l_@@_query_scheme_tl {#2} }
          {
%    \end{macrocode}
% \opt{official} 模式的韵母是音位形式，可能以 |v| 开头（|yue| 还原成 |ve|）。
% \cs{@@_pinyin:n} 只在 |l| 与 |n| 之后才把 |v| 当作 \texttt{ü}
% （见 \cs{@@_replace_v:n}：|ju|、|qu|、|xu| 没有歧义，所以不必标 \texttt{ü}），
% 开头的 |v| 会排成 |u|。因此这种形式不做标调，按无调形式输出，
% 手册里说明这一限制。判断要放在“取韵母”这一支内部，
% 不能提到外层——否则 \tn{xpinyinvalue} 也会被 \opt{scheme} 影响，
% 而它返回的是完整读音，与切分方式无关（实测提到外层时
% \opt{official} 下 \tn{xpinyinvalue} 会返回 |uang| 而不是 |wáng|）。
%    \begin{macrocode}
            \str_if_eq:nnTF {#1} { yunmu }
              { \@@_query_yunmu_out:n {#2} }
              { \@@_query_tone_apply:n {#2} }
          }
      }
  }
%    \end{macrocode}
% 先可展开地取出要的那一段（仍是“字母加声调数字”的形式），
% 再统一交给 \cs{@@_query_tone_apply:n} 决定输出形式。
% 这样标调这一步只在最外层出现一次：\cs{@@_query_mark:n} 依赖
% \cs{@@_pinyin:n}，不可展开，若嵌在 \texttt{f} 类展开里会散架
% （实测会把 \pkg{expl3} 的内部命令原样排进页面）。
%    \begin{macrocode}
\cs_new:Npn \@@_query_yunmu_join:nn #1#2 { #1 #2 }
%    \end{macrocode}
% 韵母的输出还有一处限制。\cs{@@_pinyin:n} 只在 |l| 与 |n| 之后把 |v| 当作
% \texttt{ü}（见 \cs{@@_replace_v:n}），所以韵母\emph{以 |v| 开头}时标不出
% \texttt{ü}：\verb|\pinyin{v3}| 排出来是 |ǔ| 而不是 |ǚ|。
% “女”的韵母正是这种情况。与其静默给出错的字形，这里改为输出无调形式。
% \opt{official} 模式的韵母是音位形式（|yue| 还原成 |ve|），同样不标调。
%    \begin{macrocode}
\cs_new:Npn \@@_query_yunmu_out:n #1
  {
    \exp_args:Nf \@@_query_yunmu_out_aux:nn
      { \exp_args:NV \@@_query_yunmu:nn \l_@@_query_scheme_tl {#1} } {#1}
  }
\cs_new:Npn \@@_query_yunmu_out_aux:nn #1#2
  {
    \bool_lazy_or:nnTF
      { \str_if_eq_p:ee { \l_@@_query_scheme_tl } { official } }
      { \str_if_eq_p:ee { \str_range:nnn {#1} {1} {1} } { v } }
      { #1 }
      {
        \exp_args:Nf \@@_query_tone_apply:n
          { \exp_args:Nnf \@@_query_yunmu_join:nn {#1}
              { \@@_query_tone:n {#2} } }
      }
  }
%    \end{macrocode}
% \end{macro}
%
% \begin{macro}[int]{\@@_query_head:n}
% 取字母部分的首字符，供 \tn{xpinyininitial} 使用。
%    \begin{macrocode}
\cs_new:Npn \@@_query_head:n #1 { \str_range:nnn {#1} {1} {1} }
%    \end{macrocode}
% \end{macro}
%
% \begin{macro}{\xpinyinvalue,\xpinyininitial,\xpinyinshengmu,\xpinyinyunmu}
% 四个查询命令。带星号的形式返回全部读音，用逗号分隔；不带星号取首选读音，
% 与 \tn{xpinyin} 的取值一致。
%
% 这些命令\emph{不接受可选参数}。\opt{scheme} 与 \opt{tone} 用 \tn{xpinyinsetup}
% 设置，与包内其他选项一致，并随 \TeX{} 分组恢复。这样做是因为在可展开的命令里
% 无法调用 \cs{keys_set:nn}：keyval 解析本身不可展开，若把选项放进参数，
% 整条链就用不到 \tn{index} 或 \pkg{bib2gls} 的字段里去了。
%
% \opt{tone} 取 \opt{mark} 时输出带调号，这一步依赖 \cs{@@_pinyin:n}，不可展开；
% 其余三种取值下四个命令都是完全可展开的。要拿到字面文字（例如自己 \tn{edef}
% 出索引排序键）时，用 \opt{number}、\opt{none} 或 \opt{v}。
%    \begin{macrocode}
\NewExpandableDocumentCommand \xpinyinvalue { s m }
  { \@@_query_dispatch:nnn { value } {#1} {#2} }
\NewExpandableDocumentCommand \xpinyininitial { s m }
  { \@@_query_dispatch:nnn { initial } {#1} {#2} }
\NewExpandableDocumentCommand \xpinyinshengmu { s m }
  { \@@_query_dispatch:nnn { shengmu } {#1} {#2} }
\NewExpandableDocumentCommand \xpinyinyunmu { s m }
  { \@@_query_dispatch:nnn { yunmu } {#1} {#2} }
%    \end{macrocode}
% \end{macro}
%
% \begin{macro}[int]{\@@_query_dispatch:nnn}
% 分派到“取首选读音”或“取全部读音”。|#2| 是星号参数。
%    \begin{macrocode}
\cs_new:Npn \@@_query_dispatch:nnn #1#2#3
  {
%    \end{macrocode}
% 先检查参数是不是恰好一个字符。取码位用的 |`#1| 只接受单个字符记号，
% 参数为空或多于一个字符时会漏出 \TeX{} 底层错误
% （实测 \verb|\xpinyinvalue{中国}| 先报 \texttt{! Missing \string\endcsname inserted.}，
% 空参数报 \texttt{! Improper alphabetic constant.}），看不出与本宏包有关。
%
% 这里数的是\emph{字符}而不是记号：\cs{tl_count:n} 把 \verb|{中}|、
% \verb|\hz| 这类各算一个记号，检查会放它们过去，随后照旧漏出底层错误。
% \cs{str_count:n} 先把参数变成字符串再数，能一并挡住。两者都可展开。
%    \begin{macrocode}
    \int_compare:nNnTF { \str_count:n {#3} } = { 1 }
      {
        \bool_if:nTF {#2}
          { \@@_query_all:nn {#1} {#3} }
          { \@@_query_first:nn {#1} {#3} }
      }
      { \@@_query_bad_arg:n {#3} }
  }
%    \end{macrocode}
% 首选读音直接查单音字表。查询表没有载入时（未使用 \opt{query} 选项），
% 对应的控制序列不存在，这里给出明确的报错而不是静默展开成空。
%    \begin{macrocode}
\cs_new:Npn \@@_query_first:nn #1#2
  {
    \cs_if_exist:cTF { c_@@_query_ \int_eval:n { `#2 } _tl }
      {
        \exp_args:Nnf \@@_query_part:nn {#1}
          { \@@_query_get:n {#2} }
      }
      { \@@_query_missing:n {#2} }
  }
%    \end{macrocode}
% 取不到值时分两种情况：查询表根本没载入，或者这个字不在数据库里。
% 两者的处理办法不同，所以报错也要分开。
%
% 报错必须走 \cs{msg_expandable_error:nnn} 而不是 \cs{msg_error:nnn}。
%
% 查询命令的用途之一是先 \tn{edef} 出字面文字再交给别的命令（例如索引排序键），
% 而 \cs{msg_error:nnn} 是 \texttt{protected} 的：在 \tn{edef} 里它既不报错也不
% 展开，会把自己整条原样留在结果里。实测未加 \opt{query} 选项时，
% \verb|\edef\x{\xpinyinvalue{汉}}| 得到的是
% \verb|\__xpinyin_query_missing:n {汉}|，文档零错误编译，
% 错误的排序键被静默写进 \file{.idx}。
%
% \cs{msg_expandable_error:nnn} 在可展开的上下文里也会发出错误，且不留下残余。
%    \begin{macrocode}
\cs_new:Npn \@@_query_bad_arg:n #1
  { \msg_expandable_error:nnn { xpinyin } { query-bad-argument } {#1} }
\cs_new:Npn \@@_query_missing:n #1
  {
    \bool_if:NTF \g_@@_query_loaded_bool
      { \msg_expandable_error:nnn { xpinyin } { query-no-reading } {#1} }
      { \msg_expandable_error:nn  { xpinyin } { query-not-loaded } }
  }
\bool_new:N \g_@@_query_loaded_bool
%    \end{macrocode}
% 星号形式把多音字的每个读音各加工一遍，再用逗号连接。
% 没有多音字记录的字（单音字）退回首选读音，
% 使星号形式对任何汉字都有合理结果。
%    \begin{macrocode}
\cs_new:Npn \@@_query_all:nn #1#2
  {
    \cs_if_exist:cTF { c_@@_query_multiple_ \int_eval:n { `#2 } _clist }
      {
        \exp_args:Nnf \@@_query_all_aux:nn {#1}
          { \@@_query_get_multiple:n {#2} }
      }
      { \@@_query_first:nn {#1} {#2} }
  }
%    \end{macrocode}
% 星号形式把每个读音各加工一遍，再用逗号连接。
%
% 只有在结果是纯文本时才去重。\opt{tone} 为 \opt{mark} 的输出里含
% \cs{@@_pinyin:n}，它是 \texttt{protected} 的，\texttt{e} 类展开留不下文本
% （实测 \cs{tl_set:Nx} 之后仍是 \verb|\__xpinyin_pinyin:n {zhong1}|），
% 收不进逗号列表，去重也就无从进行——强行去重会把每一项都变成空
% （实测输出 \verb|,,,,|）。
%
% 因此 \opt{mark} 下逐项输出、不去重；其余三种取值下去重。
% 首字母不受 \opt{tone} 影响，总是纯文本，所以“行”的四个读音会正确合成
% |h,x| 两项。
%    \begin{macrocode}
\cs_new_protected:Npn \@@_query_all_aux:nn #1#2
  {
    \bool_lazy_and:nnTF
      { \str_if_eq_p:ee { \l_@@_query_tone_tl } { mark } }
      { ! \str_if_eq_p:nn {#1} { initial } }
      { \@@_query_all_plain:nn {#1} {#2} }
      { \@@_query_all_unique:nn {#1} {#2} }
  }
%    \end{macrocode}
% 不去重的一支：逐项排出，项间补逗号。
%    \begin{macrocode}
\cs_new_protected:Npn \@@_query_all_plain:nn #1#2
  {
    \bool_set_true:N \l_@@_query_first_bool
    \clist_map_inline:nn {#2}
      {
        \bool_if:NTF \l_@@_query_first_bool
          { \bool_set_false:N \l_@@_query_first_bool } { , }
        \@@_query_part:nn {#1} {##1}
      }
  }
\bool_new:N \l_@@_query_first_bool
%    \end{macrocode}
% 去重的一支。
%    \begin{macrocode}
\cs_new_protected:Npn \@@_query_all_unique:nn #1#2
  {
    \clist_clear:N \l_@@_query_join_clist
    \clist_map_inline:nn {#2}
      {
        \clist_put_right:Ne \l_@@_query_join_clist
          { \@@_query_part:nn {#1} {##1} }
      }
    \clist_remove_duplicates:N \l_@@_query_join_clist
    \clist_use:Nn \l_@@_query_join_clist { , }
  }
%    \end{macrocode}
% 去重用的逗号列表变量。多音字的不同读音常落在同一个首字母上（“行”的四个读音里
% 三个都是 |h|），逐项列出没有意义。用到局部变量，所以星号形式不可展开——
% 需要字面文字时用不带星号的形式。
%    \begin{macrocode}
\clist_new:N \l_@@_query_join_clist
%    \end{macrocode}
% \end{macro}
%
%    \begin{macrocode}
\cs_if_exist:NTF \ProcessKeyOptions
  { \ProcessKeyOptions [ xpinyin ] }
  {
    \RequirePackage { l3keys2e }
    \ProcessKeysOptions { xpinyin }
  }
%    \end{macrocode}
%
% 选项处理完毕后再决定要不要载入查询表。放在这里是因为 \opt{query} 的值
% 到这一步才确定。查询命令只在 Unicode 引擎下可用，所以先检查引擎。
%    \begin{macrocode}
\bool_if:NT \g_@@_query_bool
  {
%    \end{macrocode}
% 判据写成“必须是 \XeTeX{}”而不是“不能是 \pdfTeX{}”：包级门禁目前只放行
% \XeTeX{} 与 \pdfTeX{}，但日后若放行别的引擎，正向判断不会把它默认当成可用。
%    \begin{macrocode}
    \sys_if_engine_xetex:TF
      {
        \group_begin:
          \cs_set_eq:NN \XPYUQ  \xpinyin_query:nn
          \cs_set_eq:NN \XPYUQM \xpinyin_query_multiple:nn
          \file_input:n { xpinyin-query.def }
        \group_end:
        \bool_gset_true:N \g_@@_query_loaded_bool
      }
      { \msg_error:nn { xpinyin } { query-needs-unicode } }
  }
%    \end{macrocode}
%
%    \begin{macrocode}
%</package>
%    \end{macrocode}
%
% \section{\pkg{xpinyin.lua}}
%
%    \begin{macrocode}
%<*lua>
%    \end{macrocode}
%
%    \begin{macrocode}
xpinyin       = xpinyin or { }
local xpinyin = xpinyin
%    \end{macrocode}
%
% 计算时区\footnote{\url{http://lua-users.org/wiki/TimeZone}}。
%    \begin{macrocode}
xpinyin.tzoffset = "+0000"
do
  -- Compute the difference in seconds between local time and UTC.
  local function get_timezone()
    local now = os.time()
    return os.difftime(now, os.time(os.date("!*t", now)))
  end
  -- Return a timezone string in ISO 8601:2000 standard form (+hhmm or -hhmm)
  local function get_tzoffset(timezone)
    local h, m = math.modf(timezone / 3600)
    return string.format("%+.4d", 100 * h + 60 * m)
  end
  xpinyin.tzoffset = get_tzoffset(get_timezone())
end
%    \end{macrocode}
%
%    \begin{macrocode}
xpinyin = {
  uchar    = unicode.utf8.char,
  readings = { },
  fixreadings = {
%    \end{macrocode}
% 为汉字“〇”增加拼音。
%    \begin{macrocode}
    {"U+3007", "Mandarin", "líng"}
  },
%    \end{macrocode}
%
%    \begin{macrocode}
  database = {
    source  = "https://www.unicode.org/Public/UNIDATA/Unihan.zip",
    file    = "Unihan_Readings.txt",
    date    = "Date: 2014-05-09 18:17:02 GMT [JHJ]",
    version = "Unicode version: 7.0.0",
    dbfile  = "xpinyin.db",
    qdbfile = "xpinyin-query.db"
  },
  preamble = [[
%<<COMMENT
%%
%%  Do not edit this file!
%%  Created from Unihan database:
%%
%%    $file
%%    $date
%%    $version
%%
%%  by "texlua xpinyin.lua" on ]]
           .. os.date("%Y-%m-%d %X ") .. xpinyin.tzoffset
           .. "\n%%"
%COMMENT
}
%    \end{macrocode}
%
%    \begin{macrocode}
local zip_open = require("zip").open
%    \end{macrocode}
%
% 将 \file{Unihan_Readings.txt}^^A
% \footnote{\url{http://http://www.unicode.org/reports/tr38/}。} 保存到一张表里面。
%    \begin{macrocode}
function xpinyin.maketable (txt)
  local f = io.open(txt or xpinyin.database.file, "r")
  if not f then
    local source = xpinyin.database.source
    local zfilename = source:match("[^/]+$")
    local zfile = zip_open(zfilename)
    if not zfile then
      xpinyin.download(source, zfilename)
      zfile = assert(zip_open(zfilename))
    end
    f = assert(zfile:open(xpinyin.database.file))
    zfile:close()
  end
  local s, prop
  for line in f:lines() do
    s = line:explode("\t")
    if #s == 3 then
      prop = s[2]:sub(2)
      if prop == "Mandarin" or
         prop == "HanyuPinyin" or
         prop == "XHC1983" or
         prop == "HanyuPinlu" then
        xpinyin.insert(s[1], prop, s[3])
      end
    elseif line:find("Date") then
      xpinyin.database.date = line:match("^[#%s]*(.*)")
    elseif line:find("Unicode version:") then
      xpinyin.database.version = line:match("^[#%s]*(.*)")
    end
  end
  f:close()
  if xpinyin.fixreadings then
    for _, s in pairs(xpinyin.fixreadings) do
      xpinyin.insert(s[1], s[2], s[3])
    end
  end
end
%    \end{macrocode}
%
% 下载 \file{Unihan.zip}。
%    \begin{macrocode}
function xpinyin.download (source, zip)
  print("\nRetrieving Unihan Database from\n", source)
  local ok, http = pcall(require, "socket.http")
  if ok then
    local lok, ltn12 = pcall(require, "ltn12")
    if lok then
      local fh = io.open(zip, "wb")
      if fh then
        local status = http.request{
          url  = source,
          sink = ltn12.sink.file(fh)
        }
        if status then return end
      end
    end
  end
  local ret = os.execute('curl -fsSL -o "' .. zip .. '" "' .. source .. '"')
  if ret == 0 or ret == true then return end
  ret = os.execute('wget -q -O "' .. zip .. '" "' .. source .. '"')
  if ret == 0 or ret == true then return end
  error(
    "Failed to download '" .. zip .. "'. Try one of:\n"
    .. '  - Run: texlua --socket xpinyin.lua\n'
    .. '  - Install curl or wget\n'
    .. "  - Manually place " .. zip .. " in the current directory"
  )
end
%    \end{macrocode}
%
% 往拼音表中加入项目。
%    \begin{macrocode}
function xpinyin.insert (unicode, prop, value)
  local index = tonumber(unicode:match("%x+$"), 16)
  if not xpinyin.readings[index] then
    xpinyin.readings[index] = { }
  end
  xpinyin.readings[index][prop] = value
end
%    \end{macrocode}
%
% 把带调号的拼音转成“字母加声调数字”的形式，例如 |zhōng| 转成 |zhong1|、
% |nǚ| 转成 |nv3|。这一步在生成数据库时做，查询命令运行时只需查表，
% 于是可以保持完全可展开（见查询命令一节的说明）。
%
% \file{Unihan} 里的拼音用的是预组合字符，例如 |ō| 是单独一个码位 U+014D，
% 而不是 |o| 加上长音符两个码位。因此这里用一张显式的对照表，
% 而不是依赖 Unicode 正规化——\file{texlua} 没有提供正规化函数。
% 表中列出该数据库实际用到的全部 29 个带附加符号的字符，
% 每项给出“去掉附加符号后的字母”和“声调数字”。
%
% |ü| 系需要单独说明。按包内 \tn{pinyin} 的输入约定，|ü| 写作 |v|
% （手册里的 \verb|\pinyin{nv3hai2zi}|），所以 |ǚ| 的字母部分是 |v| 而不是 |u|：
% 若写成 |u|，|nǚ| 会变成 |nu3|，馈回 \tn{pinyin} 得到的是“努”而不是“女”。
%
% 轻声没有调号，声调数字为空字符串，例如 |de| 仍然是 |de|。
%    \begin{macrocode}
xpinyin.numbertable = {
  ["ā"] = {"a", "1"}, ["á"] = {"a", "2"},
  ["ǎ"] = {"a", "3"}, ["à"] = {"a", "4"},
  ["ē"] = {"e", "1"}, ["é"] = {"e", "2"},
  ["ě"] = {"e", "3"}, ["è"] = {"e", "4"},
  ["ī"] = {"i", "1"}, ["í"] = {"i", "2"},
  ["ǐ"] = {"i", "3"}, ["ì"] = {"i", "4"},
  ["ō"] = {"o", "1"}, ["ó"] = {"o", "2"},
  ["ǒ"] = {"o", "3"}, ["ò"] = {"o", "4"},
  ["ū"] = {"u", "1"}, ["ú"] = {"u", "2"},
  ["ǔ"] = {"u", "3"}, ["ù"] = {"u", "4"},
  ["ǖ"] = {"v", "1"}, ["ǘ"] = {"v", "2"},
  ["ǚ"] = {"v", "3"}, ["ǜ"] = {"v", "4"},
  ["ü"] = {"v", "" },
%    \end{macrocode}
% |ń|、|ň|、|ǹ| 和 |ḿ| 用在“嗯”一类字上，声调直接标在鼻音字母上。
%    \begin{macrocode}
  ["ń"] = {"n", "2"}, ["ň"] = {"n", "3"}, ["ǹ"] = {"n", "4"},
  ["ḿ"] = {"m", "2"}
}
%    \end{macrocode}
%
% 少数组合并没有对应的预组合字符，只能以“字母加组合符”两个码位出现。
% 本数据库里的实例是 U+5463 的读音 |m̀|：Unicode 有 |ḿ|（U+1E3F），
% 却没有 |m| 加钝音符的预组合形式。因此除上表之外，还要单独认出裸组合符。
%    \begin{macrocode}
xpinyin.combining = {
  ["\u{0304}"] = "1", ["\u{0301}"] = "2",
  ["\u{030C}"] = "3", ["\u{0300}"] = "4"
}
xpinyin.diaeresis = "\u{0308}"
function xpinyin.numbered (pinyin)
  local base, tone = { }, ""
  for _, ch in ipairs(xpinyin.chars(pinyin)) do
    local item = xpinyin.numbertable[ch]
    if item then
      base[#base + 1] = item[1]
      if item[2] ~= "" then tone = item[2] end
    elseif xpinyin.combining[ch] then
      tone = xpinyin.combining[ch]
    elseif ch == xpinyin.diaeresis then
%    \end{macrocode}
% 裸分音符只会跟在 |u| 后面，同样按 \tn{pinyin} 的约定折叠成 |v|。
%    \begin{macrocode}
      if base[#base] == "u" then base[#base] = "v" end
    else
      base[#base + 1] = ch
    end
  end
  return table.concat(base) .. tone
end
%    \end{macrocode}
%
% 把字符串按码位切成单个字符。
%    \begin{macrocode}
function xpinyin.chars (s)
  local t = { }
  for _, cp in utf8.codes(s) do t[#t + 1] = xpinyin.uchar(cp) end
  return t
end
%    \end{macrocode}
%
% 输出需要的格式文件。
%    \begin{macrocode}
function xpinyin.output (db)
  local f = assert(io.open(db or xpinyin.database.dbfile, "w"))
  local preamble = xpinyin.preamble:gsub("%$(%w+)", xpinyin.database)
  f:write(preamble, "\n")
  local hanzi, pinyin
  local mt, qt = { }, { }
  for index, pyt in xpinyin.pairsByKeys(xpinyin.readings) do
    pinyin = assert(xpinyin.grep(pyt))
    hanzi = xpinyin.uchar(index)
    f:write("\\XPYU{", hanzi, "}{", index, "}{", pinyin, "}\n")
%    \end{macrocode}
% 查询用的“字母加声调数字”形式写进另一个文件。查询表只有用到查询命令的文档
% 才需要，分开存放使 \opt{query} 选项可以按需载入，不改变现有用户的开销。
%    \begin{macrocode}
    qt[#qt + 1] = "\\XPYUQ{" .. index .. "}{"
                  .. xpinyin.numbered(pinyin) .. "}"
    pinyin = xpinyin.multiple(pyt)
    if pinyin then
      mt[#mt + 1] = "\\XPYUM{" .. hanzi .. "}{" .. index .. "}{" .. pinyin .."}"
      local nums = { }
      for _, py in ipairs(pinyin:explode(",")) do
        nums[#nums + 1] = xpinyin.numbered(py)
      end
      qt[#qt + 1] = "\\XPYUQM{" .. index .. "}{"
                    .. table.concat(nums, ",") .. "}"
    end
  end
  f:write(table.concat(mt, "\n"), "\n")
  f:close()
  local q = assert(io.open(xpinyin.database.qdbfile, "w"))
  q:write(preamble, "\n")
  q:write(table.concat(qt, "\n"), "\n")
  q:close()
end
%    \end{macrocode}
%
% 将表按照索引排序，代码来源于
% \href{http://www.lua.org/pil/19.3.html}{\textit{Programming in Lua}}。
%    \begin{macrocode}
function xpinyin.pairsByKeys (t, f)
  local a = { }
  for n in pairs(t) do a[#a + 1] = n end
  table.sort(a, f)
  local i = 0 -- iterator variable
  return function () -- iterator function
    i = i + 1
    return a[i], t[a[i]]
  end
end
%    \end{macrocode}
%
% 按照 \texttt{Mandarin}、\texttt{XHC1983}、\texttt{HanyuPinyin}
% 的顺序选择最常用的拼音。\texttt{HanyuPinlu} 的质量较差，不采用。
%    \begin{macrocode}
function xpinyin.grep (pyt)
  if pyt.Mandarin then
    return pyt.Mandarin:match("%S+"), "Mandarin"
  elseif pyt.XHC1983 then
    return pyt.XHC1983:match(":(%S+)"), "XHC1983"
  elseif pyt.HanyuPinyin then
    return pyt.HanyuPinyin:match(":([^,%s]+)"), "HanyuPinyin"
  end
end
%    \end{macrocode}
%
% 根据 \texttt{XHC1983} 和 \texttt{HanyuPinyin} 选出多音字。
%    \begin{macrocode}
function xpinyin.multiple (pyt)
  if pyt.XHC1983 then
    local s = pyt.XHC1983:explode()
    if s[2] then
      local t = { }
      for i, v in ipairs(s) do
        t[#t + 1] = v:explode(":")[2]
      end
      return xpinyin.unique(t), "XHC1983"
    end
  elseif pyt.HanyuPinyin and pyt.HanyuPinyin:find("%D,") then
    local t = { }
    for _, v in ipairs(pyt.HanyuPinyin:explode()) do
      for _, py in ipairs(v:explode(":")[2]:explode(",")) do
        t[#t + 1] = py
      end
    end
    return xpinyin.unique(t), "HanyuPinyin"
  end
end
%    \end{macrocode}
%
% 删除掉数组中的重复元素。
%    \begin{macrocode}
function xpinyin.unique (t)
  local rt = xpinyin.remove_duplicate(t)
  if #rt > 1 then
    return table.concat(rt, ",")
  end
end
function xpinyin.remove_duplicate (t)
  local ht = { }
  local nt = { }
  for i, v in ipairs(t) do
    if not ht[v] then
      nt[#nt + 1] = v
      ht[v] = true
    end
  end
  return nt
end
%    \end{macrocode}
%
%    \begin{macrocode}
xpinyin.maketable()
%    \end{macrocode}
%
%    \begin{macrocode}
xpinyin.output()
%    \end{macrocode}
%
%    \begin{macrocode}
%</lua>
%    \end{macrocode}
%
% \end{implementation}
%
% \Finale
%
\endinput
